% concatenated from girth5.tex
\documentclass[a4paper,12pt]{article}
\usepackage{xcolor}
\usepackage{xifthen}

%% begin fonts setting
\def\fontsettingup{2} 
% 1: euler script; 2: mathpazo

\usepackage{amsmath, amssymb}
\usepackage[T1]{fontenc}
\ifthenelse{\fontsettingup = 1}{ \usepackage{eulerpx, eucal, tgpagella}   }{}
\ifthenelse{\fontsettingup = 2}{ \usepackage{mathpazo, tgpagella} }{}
%% end fonts setting

%% begin config for links %%
\usepackage{hyperref}
\hypersetup{
  colorlinks=true,
  linkcolor=blue,
  citecolor=blue,
  filecolor=magenta,
  urlcolor=cyan,
}
%% end config for links %%

\usepackage{algorithm2e}
\usepackage{tikz}
\usepackage{crossreftools} % use \crtCref{} in \section{} and \subsection{}

% start config pages %%
\usepackage[a4paper,margin=1in]{geometry}
\usepackage{microtype}
\usepackage{enumitem}
\usepackage{needspace}
\setlist[itemize]{leftmargin=1.5em,itemsep=0.3em,topsep=0.3em}
\setlength{\emergencystretch}{2em}
\allowdisplaybreaks[2]
% end config pages %%

%%   begin theorem environment  %%
\usepackage{amsthm, thmtools}
\usepackage[capitalise,nameinlink]{cleveref}

% Macro to automatically loop and define shared-counter environments
\newcommand{\addsharedenv}[1]{\declaretheorem[sibling=theorem]{#1}}

\theoremstyle{plain} 
\declaretheorem{theorem} % Base counter
\forcsvlist{\addsharedenv}{observation, claim, fact, lemma, proposition, corollary}

\theoremstyle{definition} 
\forcsvlist{\addsharedenv}{definition, remark, condition, assumption, example}
%% end theorem environment  %%

%% begin customized macro %%
% euler script setting up
\ifthenelse{\fontsettingup = 1}{
  \def\*#1{\mathbf{#1}} % Use \*A for \mathbf{A}
  \def\+#1{\mathcal{#1}} % Use \+A for \mathcal{A}
  \def\-#1{\mathrm{#1}} % Use \-A for \mathrm{A}
  \def\^#1{\mathbb{#1}} % Use \^A for \mathbb{A}
  \def\!#1{\mathfrak{#1}} % Use \!A for \mathfrak{A}
}{}

% tgpagella setting up
\ifthenelse{\fontsettingup = 2}{
  \def\*#1{\boldsymbol{#1}} % Use \*A for \mathbf{A}
  \def\+#1{\mathcal{#1}} % Use \+A for \mathcal{A}
  \def\-#1{\mathrm{#1}} % Use \-A for \mathrm{A}
  \def\=#1{\mathbb{#1}} % Use \^A for \mathbb{A}
  \def\!#1{\mathfrak{#1}} % Use \!A for \mathfrak{A}
}{}

\def\fonttest{$\lambda, \*\lambda, A, \*A, \+A, \+C, \^{R}, \^C, \^F, \^I, \!A, \!R, \!I$}

% argmin & argmax
\DeclareMathOperator*{\argmax}{arg\,max}
\DeclareMathOperator*{\argmin}{arg\,min}

% probability
% \DeclareMathOperator*{\oPr}{\mathbf{Pr}}
\def\oPr{\mathbf{Pr}}
\renewcommand{\Pr}[2][]{ \ifthenelse{\isempty{#1}}
  {\oPr\left[#2\right]}
  {\oPr_{#1}\left[#2\right]} } % Use \Pr[a]{b} for \mathbf{Pr}_a[b], \Pr{b} for  \mathbf{Pr}[b]

% expectation: relies on the package "dsfont"
% \DeclareMathOperator*{\oE}{\mathbb{E}}
\def\oE{\mathbb{E}}
\newcommand{\E}[2][]{ \ifthenelse{\isempty{#1}}
  {\oE\left[#2\right]}
  {\oE_{#1}\left[#2\right]} }

\usepackage{xparse}

% variance
% \DeclareMathOperator*{\oVar}{\mathbf{Var}}
\def\oVar{\mathbf{Var}}
% \newcommand{\Var}[2][]{ \ifthenelse{\isempty{#1}}
% {\oVar\left[#2\right]}
% {\oVar_{#1}\left[#2\right]} }
\NewDocumentCommand{\Var}{ O{} O{} m }{
  \ifthenelse{\isempty{#1}} {
    \ifthenelse{\isempty{#2}} {
      \oVar\left[#3\right]
    } {
      \oVar^{#2}\left[#3\right]
    }
  } {
    \ifthenelse{\isempty{#2}} {
      \oVar_{#1}\left[#3\right]
    } {
      \oVar_{#1}^{#2}\left[#3\right]
    }
  }
}

% entropy
% \DeclareMathOperator*{\oEnt}{\mathbf{Ent}}
\def\oEnt{\mathbf{Ent}}
% \newcommand{\Ent}[2][]{ \ifthenelse{\isempty{#1}}
% {\oEnt\left[#2\right]}
% {\oEnt_{#1}\left[#2\right]} }
\NewDocumentCommand{\Ent}{ O{} O{} m }{
  \ifthenelse{\isempty{#1}} {
    \ifthenelse{\isempty{#2}} {
      \oEnt\left[#3\right]
    } {
      \oEnt^{#2}\left[#3\right]
    }
  } {
    \ifthenelse{\isempty{#2}} {
      \oEnt_{#1}\left[#3\right]
    } {
      \oEnt_{#1}^{#2}\left[#3\right]
    }
  }
}
%% end customized macro %%

% permanent used
\newcommand{\DKL}[2]{\-D_{\-{KL}}\left(#1 \parallel #2\right)}
\newcommand{\DTV}[2]{\-D_{\mathrm{TV}}\left({#1},{#2}\right)}
\newcommand{\dist}{\mathrm{dist}}
\newcommand{\e}{\mathrm{e}}
\renewcommand{\epsilon}{\varepsilon}
\renewcommand{\emptyset}{\varnothing}
\newcommand{\norm}[1]{\left\Vert#1\right\Vert}
\newcommand{\set}[1]{\left\{#1\right\}}
\newcommand{\tuple}[1]{\left(#1\right)}
\newcommand{\eps}{\varepsilon}
\newcommand{\inner}[2]{\left\langle #1,#2\right\rangle}
\newcommand{\tp}{\tuple}
\newcommand{\ol}{\overline}
\newcommand{\abs}[1]{\left\vert#1\right\vert}
\newcommand{\ctp}[1]{\left\lceil#1\right\rceil}
\newcommand{\ftp}[1]{\left\lfloor#1\right\rfloor}
\newcommand{\Ind}[1]{\-{Ind}_{#1}}

\newcommand{\Ker}{\mathsf{Ker}\;}
\newcommand{\Ran}{\mathsf{Ran}\;}
%\newcommand{\indep}{\mathrel{\perp\!\!\!\perp}}
\newcommand{\indep}{\mathrel{\perp\mkern-10mu\perp}}
\newcommand{\todo}[1]{{\color{red}TODO: #1}}


%\title{Rapid Mixing for Multi-Spin Systems \\ on Girth-Five Graphs}
\title{A Spectral Local-to-Global Principle for Spin Systems on Graphs with Girth At Least Five}
%\title{Rapid Mixing for Spin Systems on Girth-Five Graphs via a Spectral Local-to-Global Principle}
\author{Xiaoyu Chen\thanks{Email: \texttt{xiaoyu@mit.edu}, 
Massachusetts Institute of Technology. Supported by the NSF CAREER grant CCF-2443045, and the Reed Fund at MIT.} \and Kuikui Liu\thanks{Email: \texttt{liukui@mit.edu}, 
Massachusetts Institute of Technology. Supported by the NSF CAREER grant CCF-2443045, and the Reed Fund at MIT.}}
\date{}

\begin{document}
\maketitle

\begin{abstract}


It is proved that, for every $\delta \in (0,1)$, the Glauber dynamics for
the uniform distribution on proper $q$-colorings is rapidly mixing when
$q \geq (1+\delta)\Delta$ and the underlying graph has girth at least $5$
and maximum degree $\Delta = \Omega_\delta(1)$. 
This result also extends to general multi-spin systems satisfying a \emph{local spectral contraction} condition, including the anti-ferromagnetic Potts model with $q\geq (1+\delta)(1-\beta)\Delta$.
These results are achieved by a new spectral local-to-global principle on graphs with girth at least five for general multi-spin systems, and a novel Fourier analysis for Glauber dynamics on a star.
The main ideas behind all the proofs were developed through several rounds of interaction with GPT-5.6 Sol Ultra.
\end{abstract}

%\setcounter{tocdepth}{1}
\clearpage
\tableofcontents
\clearpage

\section{Introduction}

%\paragraph{Coloring.}
Sampling from the uniform distribution $\mu$ on proper $q$-colorings is
a fundamental problem in randomized algorithms, probability, and
statistical physics.  Despite decades of progress, efficient sampling
when $q$ is close to the maximum degree $\Delta$ remains a major open problem. One of the simplest and most well-studied samplers is the \emph{(heat-bath) Glauber dynamics}, which proceeds by iterating the following two-step procedure: choose a
vertex uniformly and recolor it uniformly from the colors absent from its
neighborhood. This Markov chain is reversible with stationary distribution $\mu$, and ergodic when $q \geq \Delta + 2$.

% Let $P_{\mathrm{GD}}$ denote its transition operator.  For two probability
% distributions $\nu_1$ and $\nu_2$ on $\Omega$, their total variation
% distance is
% \[
%   \DTV{\nu_1}{\nu_2}
%   :=\frac12\sum_{\sigma\in\Omega}
%     |\nu_1(\sigma)-\nu_2(\sigma)|.
% \]
% The chain is \emph{irreducible} if every state can reach every other state
% with positive probability in some number of steps.  When the chain is
% irreducible, its mixing time, for $\varepsilon\in(0,1)$, is
% \[
%   t_{\mathrm{mix}}(\varepsilon)
%   :=\min\left\{t\in\mathbb Z_{\geq0}:
%     \max_{\sigma\in\Omega}
%     \DTV{P_{\mathrm{GD}}^t(\sigma,\cdot)}{\mu}
%     \leq\varepsilon\right\}.
% \]
% We say that the chain $P$ is \emph{rapidly mixing} if
% $t_{\mathrm{mix}}(\varepsilon)$ is bounded by a polynomial in
% $n$, $\log q$, and $\log(1/\varepsilon)$.

% If $1=\lambda_1>\lambda_2\geq\cdots\geq0$ are its eigenvalues, then its
% \emph{spectral gap} is $1-\lambda_2$ and the reciprocal of the spectral gap
% is the relaxation time.  

%It is useful to distinguish this chain from two other samplers that appear
%below.  One step of the \emph{Metropolis chain} chooses a uniformly random
%vertex, proposes a uniformly random color from all $q$ colors, and accepts
%the proposal precisely when it does not create a monochromatic edge;
%otherwise the coloring is unchanged.  The \emph{flip dynamics} is a
%non-single-site chain that chooses a vertex and a second color, finds the
%maximal bichromatic component containing that vertex, and exchanges the two
%colors on this component with a probability determined by its size
%\cite{Vigoda2000}.

%\subsection{Previous results}

%The \emph{coloring reconfiguration graph} has the proper colorings as its
%vertices, with two colorings adjacent when they differ at exactly one
%vertex.
A long-standing folklore conjecture states that
Glauber dynamics mixes rapidly throughout the range $q \geq \Delta + 2$ \cite{Jerrum1995}. For general graphs of maximum degree $\Delta$, Jerrum proved optimal
$O(n\log n)$ mixing for Glauber dynamics when $q>2\Delta$ \cite{Jerrum1995}. The current state-of-the-art result is due to Carlson--Vigoda \cite{CarlsonVigoda2025}, who established rapid mixing of a Markov chain known as the flip dynamics when $q > C\Delta$ for $C \approx 1.809$, building upon prior advances by Vigoda \cite{Vigoda2000} and Chen--Delcourt--Moitra--Perarnau--Postle \cite{ChenEtAl2019}.

Stronger results are known when the graph is locally sparse.  Let
$\alpha^*\simeq1.763$ be the positive solution of
$\alpha^*=\exp(1/\alpha^*)$.  Dyer and Frieze, Hayes, and Molloy developed
burn-in and local-uniformity methods for large-girth graphs
\cite{DyerFrieze2003,Hayes2003,Molloy2004}.  Dyer, Frieze, Hayes, and Vigoda
subsequently proved optimal $O(n\log n)$ mixing of the Glauber dynamics,
for every fixed $\varepsilon>0$ and all sufficiently large $\Delta$, when either
$q>(\alpha^*+\varepsilon)\Delta$ and the girth is at least five, or
$q>(1.489+\varepsilon)\Delta$ and the girth is at least seven
\cite{DFHV2013}.  Spectral-independence arguments later reached the
$\alpha^*$ threshold on all triangle-free graphs
\cite{CGSV2021,FGYZ2022}; subsequent work obtained an almost-optimal
spectral-gap bound \cite{JPV2022} and then an optimal order-$1/n$ spectral
gap \cite{ChenFeng2025} in this regime.

The regime closest to the conjectured threshold has so far required
additional girth.  Hayes and Vigoda introduced a non-Markovian coupling
targeting the range $q\geq(1+\delta)\Delta$ at girth at least eleven in the
large-degree setting $\Delta=\Omega(\log n)$
\cite{HayesVigoda2003}.  Jain, Mizgerd, and Vigoda recently resolved the
technical obstacles in this framework and extended it to all sufficiently
large constant degrees, proving optimal $O_\delta(n\log n)$ mixing of the
Metropolis chain \cite{JMV2026}.  In a different direction, Chen, Liu, Mani,
and Moitra
proved optimal Glauber mixing for fixed $\Delta$ and $q\geq\Delta+3$ when
the girth is sufficiently large as a function of $\Delta$
\cite{CLMM2023}.
% Our result reaches the near-degree range $q\geq(1+\delta)\Delta$ already at girth five, for the relaxation time of the heat-bath Glauber dynamics.

%\subsection{Random proper coloring and Potts model}

%\paragraph{Coloring.} We show rapid mixing of Glauber dynamics for sampling random proper $q$-coloring when $q$ is close to max degree $\Delta$ as long as 
\begin{theorem}%[Girth-five coloring]
\label{thm:main}
Fix $\delta\in(0,1)$.  Let $G=(V,E)$ be a finite simple graph on $n$
vertices with maximum degree $\Delta$ and girth at least $5$, and let $q$
be an integer.  If $\Delta$ is sufficiently large in terms of $\delta$ and
\[
  q\geq (1+\delta)\Delta,
\]
then the Glauber dynamics on proper vertex $q$-colorings is irreducible,
has spectral gap
$\Omega_\delta(1/n)$, and satisfies
\[
  t_{\mathrm{mix}}(\varepsilon)
  =O_\delta\left(n^2\log q+n\log\frac1\varepsilon\right),
  \qquad \varepsilon\in(0,1).
\]
\end{theorem}

%The proof of \Cref{thm:main} is given in \Cref{sec:coloring}.

Compared with the recent result of \cite{JMV2026}, our result improves the girth requirement from $11$ to $5$.
%Previously, there was no known \textsf{FPRAS} for random proper coloring in this regime.
%The only remaining cases are triangle-free graphs and general graphs.
\begin{remark}
%    There is usually an exponential dependency of mixing time on the slack $\delta$ for proofs based on spectral independence.
%    However, the dependency of the gap $\delta$ in \Cref{thm:main} is polynomial.
One advantage of our approach compared to those based on spectral independence (see e.g. \cite{CLMM2023}) is that the dependence of our mixing time on $1/\delta$ is polynomial, not exponential.

If the maximum degree is also bounded in terms of $\delta$, that is,
$\Delta=O_\delta(1)$, then combining our result with the universality of
spectral independence~\cite{AnariEtAl2024} yields an $O_\delta(1)$ bound on
spectral independence and, consequently, the optimal $O_\delta(n\log n)$
mixing time~\cite{ChenLiuVigoda2021}.
However, this will introduce an exponential dependence on $\delta$ for the mixing time.
\end{remark}

\paragraph{Anti-ferromagnetic Potts model.}
The techniques also generalize to the anti-ferromagnetic Potts model.
Let $\beta \in [0,1]$ be the edge-activity and $q \geq 2$ be an integer.
The anti-ferromagnetic Potts model defines a Gibbs distribution $\mu$ on $G = (V,E)$ as
\begin{align*}
  \mu(\sigma) \propto \beta^{\abs{\set{uv\in E \mid \sigma_u = \sigma_v}}}, \quad \sigma \in [q]^V.
\end{align*}
In particular, when $\beta = 1$, $\mu$ becomes a product distribution; and when $\beta = 0$, $\mu$ is the uniform distribution on proper $q$-colorings.
The following generalization of \Cref{thm:main} is proved for the Potts model.

\begin{theorem}%[Girth-five anti-ferromagnetic Potts]
\label{thm:Potts}
Fix $\delta\in(0,1)$.  Let $G=(V,E)$ be a finite simple graph on $n$
vertices with maximum degree $\Delta$ and girth at least $5$, let
$\beta\in[0,1]$, and let $q$ be an integer.  If $\Delta$ is
sufficiently large in terms of $\delta$ and
\begin{align*}
  q\geq(1+\delta)(1-\beta)\Delta,
\end{align*}
then the Glauber dynamics for the anti-ferromagnetic $q$-state Potts
model with interaction parameter $\beta$ is irreducible, has spectral
gap $\Omega_\delta(1/n)$, and satisfies
\begin{align} \label{eq:spectral-gap-mixing-bound}
  t_{\mathrm{mix}}(\varepsilon)
  =
  O_\delta\left(
    n^2\log q+n\log\frac1\varepsilon
  \right),
  \qquad
  \varepsilon\in(0,1),
\end{align}
when $\beta=0$.  If $0<\beta\leq1$, then
\begin{align*}
  t_{\mathrm{mix}}(\varepsilon)
  =
  O_\delta\left(
    n^2\log q+n^2\Delta\log\frac1\beta
    +n\log\frac1\varepsilon
  \right),
  \qquad
  \varepsilon\in(0,1).
\end{align*}
\end{theorem}

We defer the proof of \Cref{thm:Potts} to \Cref{sec:Potts}.
%\Cref{thm:Potts} follows exactly the same proof blueprint as \Cref{thm:main}.
%The only difference part is the analysis of spectral gap of Glauber dynamics on stars, which needs more efforts on calculation.
Note that \Cref{thm:main} is a special case of \Cref{thm:Potts} and \Cref{thm:Potts} falls into a general framework that we will introduce below.

\subsection{General multi-spin systems}
Let $G = (V,E)$ be a graph with maximum degree $\Delta$.
Let $q \geq 2$ be an integer.
Let $A \in [0,1]^{q\times q}$ be a symmetric matrix.
Let $\*\lambda \in \mathbb{R}_{\geq 0}^q$.
Assume that at least one configuration in $[q]^V$ has positive weight.
The Gibbs distribution $\mu_G$ with \emph{interaction matrix} $A$ and \emph{external field} $\*\lambda$ on $G$ is defined by
\begin{align*}
  \mu_G(\sigma) \propto \prod_{uv \in E} A_{\sigma_u,\sigma_v} \prod_{v\in V} \lambda_{\sigma_v}, \quad \forall \sigma \in [q]^V.
\end{align*}

\begin{condition}\label{cond-local-spectral-contraction}
  Let $\epsilon \in (0,1]$ and $\alpha \geq 1$.
  We say $\mu_G$ has \emph{local $(\alpha,\epsilon)$-spectral contraction} if, with $B := \*1\*1^\intercal - A$, the following conditions hold for every $v \in V$ and every feasible pinning $\tau \in [q]^{V\setminus N[v]}$:
  \begin{itemize}
  \item it has \emph{$\alpha$-bounded marginals}:
    \begin{align*}
      \norm{B\mu^\tau_{G,v}}_{\infty} \leq \alpha/\Delta,
      \quad \text{and} \quad \norm{B\mu^\tau_{G-v,i}}_\infty \leq \alpha/\Delta, \quad \forall i \in N(v);
    \end{align*}
  \item it has \emph{$\epsilon$-contraction}:
    \begin{align*}
        \norm{
        \-{diag}\tp{\sum_{i\in N(v)} \mu^\tau_{G-v,i}}^{1/2}
        B
        \,\-{diag}\tp{\mu^\tau_{G,v}}^{1/2}
        }_{2}^2 \leq \frac{1-\epsilon}{\Delta},
    \end{align*}
  \end{itemize}
  where we treat marginal probabilities $\mu^\tau_{G,v}$ and $\mu^\tau_{G-v,i}$ as vectors in $\mathbb{R}^q$ and all norms above are  defined in Euclidean space.
\end{condition}

\begin{theorem}\label{thm:general}
  Suppose $G$ is an $n$-vertex simple graph with maximum degree $\Delta$ and girth at least $5$.
  If $\mu_G$ has local $(\alpha,\epsilon)$-spectral contraction and $\Delta = \Omega_{\alpha,\epsilon}(1) \geq 2\alpha$ and the Glauber dynamics on $\mu_G$ is irreducible, then the Glauber dynamics on $\mu_G$ has spectral gap $\Omega_{\alpha,\epsilon}(1/n)$.
\end{theorem}

The proof of \Cref{thm:general} is deferred to \Cref{sec:general}.
To obtain \Cref{thm:general}, we introduce a local-to-global theorem (\Cref{thm:local-to-global}) via the Bochner-type identity (\Cref{prop:Bochner-identity}), which reduces the spectral gap of Glauber dynamics on the whole graph (global) to the spectral gap of Glauber dynamics on stars (local).
We then analyze the spectral gap of Glauber dynamics on stars via a sophisticated Fourier analysis (\Cref{poincare-via-proj-norms}).

\begin{remark}
  The abstract framework (\Cref{thm:local-to-global} and \Cref{poincare-via-proj-norms}) is sufficiently general to also apply to antiferromagnetic
  two-spin systems, including the hard-core and Ising models, in a fixed-gap
  $\Delta$-ary uniqueness regime.  Earlier rapid-mixing results used Weitz's
  self-avoiding-walk tree~\cite{Weitz2006} to reach the corresponding
  uniqueness regime for the $(\Delta-1)$-ary tree recursion; see, for example,
  \cite{ALO2020,ChenLiuVigodaContraction2021,ChenFengYinZhang2024,ChenChenYinZhang2025}.  For fixed positive
  uniqueness slack, the distinction between the $\Delta$-ary and
  $(\Delta-1)$-ary regimes is asymptotically negligible as $\Delta\to\infty$.
  It may also be of independent interest that our method entirely avoids the
  self-avoiding-walk tree and thereby sidesteps the obstacles to extending
  Weitz-style reductions to multispin systems identified
  by~\cite{LiuManiPernice2024}.
\end{remark}

%\subsection{Proof Overview}

\subsection{Discussions about experiments with AI}
This work is mainly inspired by the recent work of G\"obel, Jenssen,
Michelen, Pappik, Perkins, and Schiller~\cite{GoebelEtAl2026}, which gives a
direct Bochner--Bakry--\'Emery proof of the spectral-gap criterion of Chen,
Chen, Chen, Yin, and Zhang~\cite{ChenEtAl2025RandomRegular}.  The original
proof of this criterion proceeds through a trickle-down theorem for field
dynamics~\cite{ChenEtAl2025RandomRegular}.

We first asked GPT-5.6 Sol Ultra to use the Bochner identity to simplify the
proofs of several results obtained by trickle-down methods, including
Oppenheim's trickle-down theorem~\cite{Oppenheim2018,AlevLau2020}, spectral independence via the pairwise
spectral-influence matrix~\cite{LeakeOveisGharan2025}, and the matrix
trickle-down theorem for sampling
edge-colorings~\cite{ALG2021}\footnote{We omit the Bochner-type proofs of
these results because these results are known before and not related to the main result of this paper.}.

With these proofs, we observe that the Bochner identity is useful for
reducing a global spectral-gap estimate to estimates on local gadgets:
\begin{itemize}
\item for Oppenheim's trickle-down theorem, the local gadgets are
  codimension-$2$ links;
\item for spectral independence via pairwise spectral influence, the local
  gadgets are two-part links;
\item for edge-coloring, the local gadgets are endpoint stars, which are
  cliques in the corresponding line graph.
\end{itemize}
These experiments also suggest that GPT is particularly effective at
analyzing such local gadgets.

It was then natural to ask whether this method works for proper vertex
colorings of graphs of girth at least $5$, where stars are again natural
local gadgets.
We brought this question to GPT-5.6 Sol Ultra, which proposed an affirmative
proof strategy.  The human authors then verified the argument, generalized
the proof with further assistance from GPT, and streamlined the paper.
The authors also use GPT-5.6 to improve the exposition and
identify typographical errors. The author takes full responsibility for the contents of this paper.

\begin{remark}[Discussion of the current obstruction]
  The existence of spanning $4$-cycles breaks both the local-to-global argument in \Cref{sec:local-to-global} and the local analysis in \Cref{sec:fourier-analysis-on-star}.
  We asked GPT-5.6 Sol Ultra to work on triangle-free graphs with the Bochner identity, but it did not find a proof after $\geq 20$ hours.

  In fact, spanning $4$-cycles seem to be the only remaining hard gadgets for coloring.
  Given the present manuscript, GPT-5.6 Sol Ultra claims that it has found a proof that generalizes \Cref{thm:main} to all graphs without spanning $4$-cycles.
  The proof follows a high-level plan similar to the one presented in this manuscript. For the local expansion, it uses the fact that, without spanning $4$-cycles, the induced subgraph on the neighbors of any vertex is a matching.
  However, this makes the analysis of the local gadget significantly more sophisticated.
  We omit this result from this manuscript since we did not find any new ideas in it.
\end{remark}

%Current bottleneck of girth requirement of \Cref{thm:main} comes from the trickle-down theorem (\Cref{thm:local-to-global}).
%It is unclear how to remove the assumption from current proof.
%It would be interesting if a similar trickle-down theorem can be proved for the tringle-free graphs.
%We leave this as a open problem.

%The order $1/n$ of the spectral gap is optimal for a single-site chain.
%The mixing bound in \Cref{thm:main} follows from the standard worst-start
%spectral-gap estimate.  This estimate loses a factor of $n$ compared with
%the optimal $O_\delta(n\log n)$ mixing time available at girth eleven.
%
%\paragraph{Proof overview.}
%The proof separates a model-independent local-to-global argument from a
%model-specific analysis on a star.  The global step is an exact Bochner
%decomposition.  After expanding the square of the Glauber Laplacian, more
%distant update pairs contribute nonnegative squares, while girth five lets
%us assign every distance-two pair to a unique closed star.  Consequently,
%a suitable weighted spectral-gap estimate on every conditioned star
%implies a global spectral-gap estimate; this is
%\Cref{thm:local-to-global}.
%
%It remains to analyze the weighted dynamics on a star.  Conditional on the
%center color, the leaves are independent.  Their conditional Hoeffding
%expansion separates degree zero, degree one, and higher-degree leaf
%components.  Writing the weighted Laplacian in these three blocks and
%eliminating the degree-zero and higher-degree blocks by Schur complements
%reduces the local Poincar{\'e} inequality to two operator norms; see
%\Cref{poincare-via-proj-norms}.  For colorings, the conditioned star admits
%an exact tree recursion.  A second-moment estimate controls the center
%channel, while contraction of the linearized tree recursion controls its
%degree-one part.  These estimates verify the local hypothesis and complete
%the proof of \Cref{thm:main}.

\subsection{Related work}

\paragraph{Rapid mixing for the anti-ferromagnetic Potts model.}
The zero-temperature case $\beta=0$ is the coloring model surveyed above.
At positive temperature, early work used Dobrushin-type and spatial-mixing
arguments to establish rapid Glauber mixing, with subsequent refinements
for bounded-degree graphs and lattices
\cite{SalasSokal1997,GoldbergEtAl2006Potts}.  Later work developed
correlation-decay and zero-free approaches to approximate counting
\cite{YinZhang2015SpatialPotts,BencsBerrekkalRegts2026Counting,
LiuSinclairSrivastava2025Potts}, clarified uniqueness and spatial mixing on
trees \cite{GalanisGoldbergYang2018,deBoerBuysRegts2023,BencsEtAl2023Potts,
CLMM2023,BencsBerrekkalRegts2025}, and obtained polynomial-time sampling on
random regular and sparse random graphs
\cite{BlancaEtAl2020Potts,YinZhang2016Potts,Efthymiou2022Potts}.
Complementary hardness and slow-mixing results describe the obstructions
below the natural scale $q=(1-\beta)\Delta$ and at low temperature
\cite{GalanisStefankovicVigoda2015,LiLiuYang2026Potts}.

\paragraph{Trickle-down methods.}
Beginning with Garland's method, local-to-global spectral arguments
propagate expansion from links to high-dimensional walks; this lineage now
includes local spectral expansion, spectral and entropic independence, and
scalar, matrix, partite, poset, localization, and $\mathcal C$-Lorentzian
trickle-down theorems
\cite{Garland1973,BallmannSwiatkowski1997,Zuk2003,
ParzanchevskiRosenthal2017,KaufmanMass2017,DinurKaufman2017,
Parzanchevski2017,Oppenheim2018,KaufmanOppenheim2020,Oppenheim2020,
AnariEtAl2019,CryanGuoMousa2021,AlevLau2020,ALO2020,GuoMousa2020,
FGYZ2022,JPV2022,ChenLiuVigoda2021,ALG2021,AnariEtAl2022,
KaufmanTessler2026,
AbdolazimiOveisGharan2023,GotlibKaufman2023,AnariKoehlerVuong2024, ChenEtAl2025RandomRegular,
LeakeLindbergOveisGharan2025,LeakeOveisGharan2025}.  Our local-to-global
theorem follows the same philosophy, but it uses an exact operator expansion
and the girth-five geometry to collect every interacting pair of updates
inside a closed star.

\paragraph{Bochner-type inequalities.}
Beginning with Bochner's identity and Bakry--\'Emery theory, comparisons
between a Dirichlet form and its iterated form have been used to derive
Poincar\'e and entropy-decay inequalities
\cite{Bochner1946,BakryEmery1985,BGL2014}.  Discrete versions were
developed for finite Markov chains and interacting-particle systems
\cite{BCDP2006,CaputoPosta2007,Caputo2008,CDP2009,DaiPraPosta2013,
Maas2011,ErbarMaas2012,Mielke2013,EMT2015,FathiMaas2016,EHMT2017,
ErbarFathi2018}, with more recent extensions involving Beckner and
curvature-dimension inequalities and coupling-based curvature bounds
\cite{JungelYue2017,WeberZacher2021,Pedrotti2025}.  Related
squared-generator expansions have also been applied to spectral-gap
estimates for hard-core and anti-ferromagnetic two-spin models
\cite{KondratievKunaOhlerich2013,GoebelEtAl2026,
GuoZhang2026Planar}.  Our Bochner identity
follows this line, using the girth-five geometry to localize pairwise
update terms on closed stars.

\section{Preliminaries}
\subsection{Notation}
\label{subsec:notation}
For a nonnegative integer $m$, write $[m]:=\{1,\ldots,m\}$, with
$[0]:=\varnothing$.  Let
$G=(V,E)$ be a finite simple graph.  We use
\begin{align*}
  n&:=|V|,
  &N(v)&:=\{u\in V:\{u,v\}\in E\},
  &N[v]&:=N(v)\cup\{v\},\\
  \deg_G(v)&:=|N(v)|,
  &\Delta&:=\max_{v\in V}\deg_G(v).
\end{align*}
Thus $N(v)$ and $N[v]$ are respectively the open and closed neighborhoods
of $v$.  We write $\dist_G(u,v)$ for graph distance, and the \emph{girth}
of $G$ is the length of its shortest cycle, with girth $\infty$ for a
forest.

A \emph{spin system} on $G$ with $q$ spins, where $q$ is a positive integer, is a
probability distribution $\mu$ whose support is
$\Omega\subseteq[q]^V$, where $[q]^V$ is the set of maps from $V$ to
$[q]$.  We write
$X=(X_v)_{v\in V}\sim\mu$ for a random configuration.  For
$B\subseteq V$, write $X_B:=(X_v)_{v\in B}$ and, for a deterministic
configuration $x$, write $x_B:=(x_v)_{v\in B}$.

For $B\subseteq V$, a \emph{pinning outside $B$} is an assignment
$\tau\in[q]^{V\setminus B}$.  It is \emph{feasible} if
$\Pr[\mu]{X_{V\setminus B}=\tau}>0$.  For a feasible pinning, we write
$\mu_B^\tau$ for the conditional law of $X_B$ given
$X_{V\setminus B}=\tau$.

For any probability distribution $\nu$ whose finite support is $\Omega$, we
write $\Pr[\nu]{A}$, $\E[\nu]{Y}$, and
$\operatorname{Var}_\nu(Y)$ for the probability of an event $A$ and the
expectation and variance of a random variable $Y$ under $\nu$; in
particular,
\begin{align*}
  \operatorname{Var}_\nu(Y)
  :=\E[\nu]{\bigl(Y-\E[\nu]{Y}\bigr)^2}.
\end{align*}
Conditional variance means variance under the indicated conditional law.
For functions on $\Omega$, define
\begin{align*}
  L^2(\nu)&:=\bigl(\mathbb R^\Omega,\inner{\cdot}{\cdot}_\nu\bigr),
  &\inner{f}{g}_\nu&:=\sum_{x\in\Omega}\nu(x)f(x)g(x).
\end{align*}
We write $\*1_A$ for the indicator of an event $A$ and $\*1$ for the
constant-one function.  The centered subspace is
\begin{align*}
  \*1^\perp
  :=\{f\in L^2(\nu):\inner{f}{\*1}_\nu=0\}
  =\{f\in L^2(\nu):\E[\nu]{f}=0\}.
\end{align*}
We use Roman letters such as $u,v,w,i,j$ for vertices and Fraktur letters
such as $\mathfrak a,\mathfrak b,\mathfrak c$ for colors.

Subscripts in $O_\delta(\cdot)$, $\Omega_\delta(\cdot)$, and
$o_\delta(\cdot)$ allow the implicit constants, thresholds, or rates of
convergence to depend on $\delta$.  In particular,
$\Delta=\Omega_\delta(1)$ means that $\Delta$ is sufficiently large in
terms of $\delta$.

For a finite set $S$, $\binom{S}{2}$ denotes the set of its two-element
subsets.

\subsection{General inner product space}
Throughout this subsection, all vector spaces are finite-dimensional.  An
inner-product space is a real vector space $H$ equipped with a symmetric
bilinear form $\inner{x}{y}_H$ such that $\inner{x}{x}_H>0$ for every
$x\neq0$.  It induces the norm $\norm{x}_H:=\sqrt{\inner{x}{x}_H}$.
For two vectors $f,g\in H$, we say $f\perp g$ if $\inner{f}{g}_H=0$.
For subspaces $K_0,K_1\subseteq H$, the notation $K_0\perp K_1$ means
that every vector in $K_0$ is orthogonal to every vector in $K_1$.  The
orthogonal complement of $K\subseteq H$ is
$K^\perp:=\{x\in H:x\perp y\text{ for every }y\in K\}$.
\begin{remark}
  For the rest of this note, once the inner-product space is clear from
  context, we use $\inner{\cdot}{\cdot}$ for its inner product and
  $\norm{\cdot}$ for the corresponding norm.
\end{remark}

For a subspace $K\subseteq H$, write $\operatorname{proj}_K$ for the
orthogonal projection onto $K$.  For a linear operator $A:H\to H$, write
$A|_K$ for its restriction to $K$, and write
\begin{align*}
  \Ker(A)&:=\{x\in H:Ax=0\},
  &\Ran(A)&:=\{Ax:x\in H\}.
\end{align*}
We write $\mathrm{Id}$ or $\mathbb I$ for the identity operator when its
space is clear.  A subscript indicates the space on which the identity
acts; for example, $\mathbb I_i$ is the identity on a subspace indexed by
$i$.

\begin{definition}
  Let $A, A^*:H\to H$ be linear operators.
  We say $A^*$ is an adjoint operator for $A$ if, for every $f,g\in H$,
  \begin{align*}
    \inner{Af}{g} = \inner{f}{A^*g}.
  \end{align*}
\end{definition}

\begin{proposition}
  Let $A:H\to H$ be a linear operator.
  Then $A^*$ exists and is unique.
\end{proposition}
The proof is standard and is omitted.

We will also use the operator analogue of positive semidefinite matrices.
\begin{definition}
  For a linear operator $A:H\to H$:
  \begin{itemize}
  \item $A$ is \emph{self-adjoint} if $A=A^*$;
  \item $A\succeq0$ means $A=A^*$ and
    $\inner{x}{Ax}_H\geq0$ for every $x\in H$;
  \item $A\succ0$ means $A=A^*$ and
    $\inner{x}{Ax}_H>0$ for every $x\neq0$;
  \item $A\succeq B$ means $A-B\succeq0$;
  \item $A\preceq B$ means $B\succeq A$;
  \item $\norm{A}:=\max_{\norm{x}_H=1}\norm{Ax}_H$ is the operator norm.
  \end{itemize}
  Thus $\succeq$ is an order on self-adjoint operators, not an entry-wise
  inequality.  All spectra below consist of eigenvalues; multiplicities are
  stated explicitly when relevant.
\end{definition}

\begin{fact}
  For every linear operator $A:H\to H$, we have
  $\norm{AA^*} = \norm{A^*A} = \norm{A}^2$.
\end{fact}

\subsection{Schur complement}
Let $H=H_0\oplus H_1$ be an orthogonal decomposition: every $x\in H$ has a
unique representation $x=x_0+x_1$, with $x_i\in H_i$ and $H_0\perp H_1$.
For a linear operator $A:H\to H$, its blocks are
\[
  A_{ij}:=\operatorname{proj}_{H_i}A|_{H_j}:H_j\longrightarrow H_i.
\]
If $A=A^*$, then
\[
  A=
  \begin{pmatrix}
    A_{00}&A_{01}\\
    A_{10}&A_{11}
  \end{pmatrix},
  \qquad
  A_{00}=A_{00}^*,\quad A_{11}=A_{11}^*,\quad A_{10}=A_{01}^*.
\]

\begin{definition}
Assume $A_{00}$ is invertible.
The \emph{Schur complement} of $A_{00}$ in $A$ is
\begin{align} \label{eq:abstract-Schur-complement}
  A/A_{00}
  :=A_{11}-A_{10}A_{00}^{-1}A_{01}.
\end{align}
\end{definition}

The following result is the main motivation behind the Schur complement.
\begin{proposition} \label{prop:Schur-complement-elimination}
  If $A=A^*$ and $A_{00}\succ0$, then
  \begin{align} \label{eq:Schur-positivity-equivalence}
    A\succeq0
    \quad\Longleftrightarrow\quad
    A/A_{00}\succeq0.
  \end{align}
\end{proposition}
\begin{proof}
Note that $A/A_{00}: H_1 \to H_1$ by definition.
Fix $z\in H_1$ and minimize
\[
  \langle h+z,A(h+z)\rangle
  \qquad\text{over }h\in H_0.
\]
Since $A_{00}\succ0$, the unique minimizer is characterized by
\[
  \operatorname{proj}_{H_0}A(h+z)
  =A_{00}h+A_{01}z=0,
\]
and hence $h=-A_{00}^{-1}A_{01}z$. Substitution gives the
variational identity
\[
  \min_{h\in H_0}\langle h+z,A(h+z)\rangle
  =\langle z,(A/A_{00})z\rangle .
\]
Because every vector in $H$ has the form $h+z$, it follows that
\begin{align*}
  A\succeq0
  &\quad\Longleftrightarrow\quad
  \min_{h\in H_0}\langle h+z,A(h+z)\rangle\geq0
  \ \text{for every }z\in H_1 \\
  &\quad\Longleftrightarrow\quad
  A/A_{00}\succeq0 . \qedhere
\end{align*}
\end{proof}

\begin{remark}
The Schur complement $A/A_{00}$ is the effective operator on $H_1$
obtained after choosing the $H_0$-component optimally.
\Cref{prop:Schur-complement-elimination}
therefore reduces positivity of an operator on $H_0\oplus H_1$ to
positivity of an operator on $H_1$. In our applications, $H_1$ has
additional structure that makes the reduced inequality easier to analyze,
even when the explicit expression for $A/A_{00}$ is more complicated.
\end{remark}

\subsection{Spectral gap, mixing time, and a Bochner identity}
Let $\Omega$ be a finite state space, let $\mu$ be a probability
distribution with full support on $\Omega$, and let $P$ be the transition
operator of a Markov chain that is reversible with respect to $\mu$.
Thus $P$ is self-adjoint on $L^2(\mu)$ and $P\*1=\*1$.  We assume that
$P\succeq0$, as will be the case for the Glauber dynamics below, and
define its discrete-time Laplacian by
\begin{align*}
  \+L:=\mathrm{Id}-P.
\end{align*}

Let
\begin{align*}
  1=\lambda_1(P)\geq\lambda_2(P)\geq\cdots\geq0
\end{align*}
be the eigenvalues of $P$, counted with multiplicity.  The
\emph{spectral gap} of $P$ is $1-\lambda_2(P)$.  It is positive exactly
when the chain is irreducible.  The ratio
\begin{align*}
  \frac{\inner{f}{\+L f}_\mu}{\inner{f}{f}_\mu}
\end{align*}
is the Rayleigh quotient of $\+L$ at $f$.  By the Rayleigh--Ritz
variational principle,
\begin{align*}
  1-\lambda_2(P)
  =\inf_{\substack{f\neq0\\f\perp\*1}}
    \frac{\inner{f}{\+L f}_\mu}{\inner{f}{f}_\mu}.
\end{align*}
Equivalently, $1-\lambda_2(P)$ is the largest constant $\gamma$ for
which the \emph{Poincar{\'e} inequality}
\begin{align*}
  \gamma\inner{f}{f}_\mu\leq\inner{f}{\+L f}_\mu,
  \qquad f\perp\*1,
\end{align*}
holds.
%
The main tool of this note is the following Bochner-type identity.
\begin{proposition} \label{prop:Bochner-identity}
  If $P$ is irreducible, then we have
  \begin{align*}
    \gamma
    :=\inf_{\substack{f\neq0\\f\perp\*1}}
      \frac{\inner{f}{\+L f}_\mu}{\inner{f}{f}_\mu}
    =\inf_{\substack{f\neq0\\f\perp\*1}}
      \frac{\inner{\+L f}{\+L f}_\mu}{\inner{f}{\+L f}_\mu}.
  \end{align*}
\end{proposition}
\begin{proof}
  We restrict to the inner-product space
  $(\*1^\perp,\inner{\cdot}{\cdot}_\mu)$.
  The left-hand side means
  $\+L\succeq\gamma\mathrm{Id}$, whereas the right-hand side means
  $\+L^2\succeq\gamma\+L$.  Since $P$ is irreducible,
  $\+L\succ0$ on $\*1^\perp$, and hence $\+L$ is invertible there.  On
  this subspace, $\+L^{1/2}$ denotes the unique positive-definite square
  root of $\+L$, and $\+L^{-1/2}$ denotes its inverse.  Conjugating
  $\+L^2\succeq\gamma\+L$ by $\+L^{-1/2}$ gives
  $\+L\succeq\gamma\mathrm{Id}$, and the converse follows by conjugating
  with $\+L^{1/2}$.
\end{proof}

For an irreducible chain, its mixing time is
\begin{align*}
  t_{\mathrm{mix}}(\varepsilon)
  :=\min\left\{t\in\mathbb Z_{\geq0}:
    \max_{x\in\Omega}\DTV{P^t(x,\cdot)}{\mu}
    \leq\varepsilon\right\},
\end{align*}
where for two probability
distributions $\nu_1$ and $\nu_2$ on $\Omega$, their total variation
distance is
\[
  \DTV{\nu_1}{\nu_2}
  :=\frac12\sum_{\sigma\in\Omega}
    |\nu_1(\sigma)-\nu_2(\sigma)|.
\]
If its spectral gap is $\gamma>0$, then, for
$0<\varepsilon\leq1/2$, positive semidefiniteness of $P$ and $L^2$
contraction give the standard worst-start bound
\cite[Theorem~12.4]{LevinPeres2017}
\begin{align}\label{eq:spgap-implies-mixing-time}
  t_{\mathrm{mix}}(\varepsilon)
  \leq\left\lceil\frac1\gamma\left(
    \log\frac1{2\varepsilon}
    +\frac12\log\frac1{\min_{x\in\Omega}\mu(x)}
  \right)\right\rceil.
\end{align}

\subsection{Glauber dynamics and its Laplacians}
We now specialize the preceding discussion to heat-bath Glauber
dynamics.
Let $V$ be a finite coordinate set with $n=|V|$, let $\mu$ be a
distribution whose support is $\Omega\subseteq[q]^V$, and let
$X\sim\mu$.  For each $v\in V$, define the single-site heat-bath operator
$P_v$ point-wise by
\begin{align} \label{def-Pi}
  (P_vf)(x)
  :=\E[\mu]{f(X)\mid
    X_{V\setminus\{v\}}=x_{V\setminus\{v\}}},
  \qquad x\in\Omega.
\end{align}
Since the state space is finite, $P_v$ can be viewed as a matrix in
$\mathbb{R}^{\Omega\times\Omega}$.  Define the single-site Laplacian by
\begin{align} \label{def-Li}
  \+L_v:=\mathrm{Id}-P_v.
\end{align}
For any $B\subseteq V$, write
\begin{align} \label{def-LB}
  \+L_B:=\sum_{v\in B}\+L_v.
\end{align}
All these operators act on the inner-product space $L^2(\mu)$ defined in
\Cref{subsec:notation}.

\begin{fact}
  The operators $P_v$ and $\+L_v$ are self-adjoint orthogonal
  projections in $L^2(\mu)$.
\end{fact}

The transition operator of the discrete-time Glauber dynamics is
\begin{align*}
  P_{\mathrm{GD}}:=\frac1n\sum_{v\in V}P_v.
\end{align*}
Define
\begin{align} \label{def-L}
  \+L:=\sum_{v\in V}\+L_v.
\end{align}
Then
\begin{align*}
  P_{\mathrm{GD}}=\mathrm{Id}-\frac1n\+L,
\end{align*}
so the Laplacian of the discrete-time Glauber dynamics is $\+L/n$.
Since $P_{\mathrm{GD}}$ is an average of orthogonal projections, it is
self-adjoint and positive semidefinite, and the preceding discussion
applies.  In particular, the spectral gap of $P_{\mathrm{GD}}$ is
$1/n$ times the smallest eigenvalue of $\+L$ on $\*1^\perp$.

In the continuous-time Glauber dynamics, every coordinate is updated at
rate one.  Its generator is $-\+L$, so $\+L$ is its Laplacian and its
spectral gap is $n$ times that of the discrete-time Glauber dynamics.

Given nonnegative update rates $a_v$, the Laplacian of the corresponding
weighted continuous-time dynamics is $\sum_{v\in V}a_v\+L_v$.  We call
this dynamics irreducible when the kernel of its Laplacian consists only
of constant functions.  The proof of \Cref{prop:Bochner-identity}
applies verbatim to every such irreducible weighted dynamics.
\section{A local-to-global principle for girth-5 spin systems} \label{sec:local-to-global}
Let $G=(V,E)$ be a graph with girth at least $5$, and use the Gibbs and
pinning conventions from \Cref{subsec:notation}.

\begin{definition}
Given a distribution on configurations on the vertex set $N[v]$, whose
induced subgraph is a star, and a parameter $0\leq\eta_v<1$, the
\emph{weighted continuous-time Glauber
dynamics with parameter $\eta_v$} has Laplacian
\begin{align*}
  \+L_v+\frac{1+\eta_v}{2}\+L_{N(v)}.
\end{align*}
Here the single-site Laplacians are formed with respect to the given
distribution: the center is updated at rate one, and each leaf is updated
at rate $(1+\eta_v)/2$.
\end{definition}

Throughout this paper, we assume that every Gibbs distribution satisfies the following conditional independence property.
\begin{definition}\label{def:cond-indep}
  A Gibbs distribution $\mu$ on $G$ is \emph{conditionally independent} if
  \begin{align*}
    X_u\indep X_v\mid X_{V\setminus\{u,v\}}
    \qquad\text{whenever }\dist_G(u,v)\geq2,
  \end{align*}
where we use $\indep$ to denote independence.
\end{definition}
%\emph{Gibbs distribution on $G$} means a spin system satisfying the
%nearest-neighbor pairwise Markov property


\begin{theorem} \label{thm:local-to-global}
Let $G=(V,E)$ be a graph with girth at least $5$, and let $\mu$ be a conditionally independent Gibbs
distribution on $G$ whose Glauber dynamics is irreducible.  Suppose that
the parameters $0\leq\eta_v<1$ satisfy the following conditions:
\begin{itemize}
\item for every $v\in V$ and every feasible pinning $\tau$ outside $N[v]$, the
weighted continuous-time Glauber dynamics for $\mu_{N[v]}^\tau$ with
parameter $\eta_v$ has spectral gap at least $(1-\eta_v)/2$;
\item there exists $\epsilon>0$ such that
$\sum_{u\in N(v)}\eta_u\leq1-\epsilon$ for every $v\in V$.
\end{itemize}
Then the continuous-time Glauber dynamics for $\mu$ has spectral gap at
least $\epsilon/2$.
\end{theorem}

\begin{proof}
Fix $f\in L^2(\mu)$.  Expanding the square, with all pairs below
unordered, gives
\begin{align*}
  \norm{\+L f}^2
  &=\sum_{v\in V}\norm{\+L_vf}^2
    +2\sum_{\{u,w\}\in\binom{V}{2}}
      \inner{\+L_u f}{\+L_w f}.
\end{align*}
If $\dist_G(u,w)\geq2$, conditional independence implies that the two
update projections commute.  Since $\+L_u$ and $\+L_w$ are commuting
orthogonal projections,
\begin{align*}
  \inner{\+L_u f}{\+L_w f}
  =\norm{\+L_u\+L_wf}^2\geq0.
\end{align*}
We therefore discard the terms with distance at least three from $\norm{\+Lf}^2$.
In other words,
\begin{align*}
  \norm{\+L f}^2
  &\geq\sum_{v\in V}\norm{\+L_vf}^2
    +2\sum_{\{u,w\}\in\binom{V}{2} \atop 1 \leq \mathrm{dist}_G(u,w) \leq 2}
    \inner{\+L_u f}{\+L_w f}.
\end{align*}

Because $G$ has no triangles, the vertices in $N(v)$ are pairwise
nonadjacent.  Thus
\begin{align*}
  \+L_{N(v)}^2-\+L_{N(v)}
  =2\sum_{\set{u,w}\in\binom{N(v)}{2}}\+L_u\+L_w\succeq0.
\end{align*}
Because $G$ has no $4$-cycles, every pair at distance two has a unique
common neighbor.
After dropping vertex pairs with distance at least three, a direct double-counting argument now gives
\begin{align*}
  \norm{\+Lf}^2
  \geq\sum_{v\in V}\underbrace{\left(
    \inner{f}{\+L_v f}
    +\inner{\+L_vf}{\+L_{N(v)}f}
    +\inner{f}{(\+L_{N(v)}^2-\+L_{N(v)})f}
  \right)}_{=:T_v}.
\end{align*}

A direct expansion shows that
\begin{align*}
  &T_v-\frac1{1+\eta_v}
  \inner{f}{\left(\+L_v+\frac{1+\eta_v}{2}\+L_{N(v)}\right)^2f}\\
  &\qquad=\inner{f}{\left(
    \frac{\eta_v}{1+\eta_v}\+L_v
    +\frac{3-\eta_v}{4}\+L_{N(v)}^2
    -\+L_{N(v)}
  \right)f}.
\end{align*}
Condition on $X_{V\setminus N[v]}$.  For each feasible pinning $\tau$, apply
the weighted spectral-gap assumption and \Cref{prop:Bochner-identity} in
$L^2(\mu_{N[v]}^\tau)$, and then average over $\tau$.  Conditioning on each such pinning,
the restriction of every global single-site Laplacian with index in $N[v]$
is the corresponding single-site Laplacian for $\mu_{N[v]}^\tau$.  The
same proof applies because the local Laplacian is a nonnegative weighted
sum of single-site Laplacians.  This gives
\begin{align*}
  &\frac1{1+\eta_v}
  \inner{f}{\left(\+L_v+\frac{1+\eta_v}{2}\+L_{N(v)}\right)^2f}\\
  &\qquad\geq\frac{1-\eta_v}{2(1+\eta_v)}
  \inner{f}{\left(\+L_v+\frac{1+\eta_v}{2}\+L_{N(v)}\right)f}.
\end{align*}
Combining the last two displays, and using
$\+L_{N(v)}^2-\+L_{N(v)}\succeq0$, yields
\begin{align*}
  T_v
  &\geq\inner{f}{\left(
    \frac12\+L_v-\frac{\eta_v}{2}\+L_{N(v)}
    +\frac{3-\eta_v}{4}
      (\+L_{N(v)}^2-\+L_{N(v)})
  \right)f}\\
  &\geq\inner{f}{\left(
    \frac12\+L_v-\frac{\eta_v}{2}\+L_{N(v)}
  \right)f}.
\end{align*}
Summing over $v$ and collecting the coefficient of each single-vertex
energy gives
\begin{align*}
  \norm{\+Lf}^2
  &\geq\frac12\sum_{v\in V}
    \left(1-\sum_{u\in N(v)}\eta_u\right)
    \inner{f}{\+L_vf}\\
  &\geq\frac\epsilon2\inner{f}{\+Lf}.
\end{align*}
The Bochner identity and irreducibility finish the proof.
\end{proof}

\section{A Fourier analysis for continuous-time Glauber dynamics on a star} \label{sec:fourier-analysis-on-star}
Let $G = (V, E)$ be a star.
For convenience, let $v$ be the center of the star and let
$N(v)=\set{1,\dots,d}$.
Let $\nu$ be a distribution on $[q]^V$, and let $X\sim\nu$.  Assume that,
conditional on $X_v$, the leaves $(X_i)_{i\in N(v)}$ are independent.
Let $P_u$ and $\+L_u$ denote the heat-bath projection
and single-site Laplacian in $L^2(\nu)$ for every $u\in V$, as defined in
\Cref{def-Pi,def-Li}.
In this section, we will work 
with the inner-product space
\begin{align*}
  (\*1^\perp, \inner{\cdot}{\cdot}_\nu).
\end{align*}


We give a general method to bound the spectral gap of
weighted continuous-time Glauber dynamics with parameter $\eta$.
Conditional independence of the leaves implies that their heat-bath
projections and single-site Laplacians commute.  Hence the order of the
products in the following definition is immaterial.

\begin{definition} \label{def:Hoeffding-expand}
Let $f \in \*1^\perp$ be a function.
Its \emph{orthogonal Hoeffding expansion} is
\begin{align*}
  f = \sum_{\Lambda \subseteq [d]} f_\Lambda, \quad \text{where} \quad f_\Lambda := \tp{\prod_{i\in\Lambda}\+L_i}\tp{\prod_{i\in[d]\setminus\Lambda}P_i}f.
\end{align*}
Here an empty product of operators is the identity.
We call $|\Lambda|$ the \emph{degree} of the component $f_\Lambda$.
\end{definition}

\begin{remark}
    The decomposition in \Cref{def:Hoeffding-expand} is well-defined since
    \begin{align*}
        \mathbb{I} = \prod_i \mathbb{I} = \prod_i\tp{\+L_i + P_i} = \sum_{\Lambda \subseteq [d]} \tp{\prod_{i\in\Lambda}\+L_i}\tp{\prod_{i\in[d]\setminus\Lambda}P_i},
    \end{align*}
    where in the last equation, we use the fact that, for all $i\neq j$, conditional independence of the leaves given the center implies that
  $\+L_i,P_i,\+L_j,P_j$ all commute.
\end{remark}

We collect some properties of the orthogonal Hoeffding expansion.
\begin{proposition} \label{prop:Hoeffding-expand}
  The following hold:
  \begin{itemize}
  \item for all distinct $A,B\subseteq[d]$, we have $f_A\perp f_B$;
  \item $\+L_i f_\Lambda = \*1_{\set{i\in\Lambda}}f_\Lambda$ and
    $\+L_{N(v)} f_\Lambda = \abs{\Lambda} f_\Lambda$.
  \end{itemize}
\end{proposition}
\begin{proof}
  For all $i\in[d]$, we have $\+L_iP_i=P_i-P_i^2=0$, since $P_i$ is
  an orthogonal projection.  Also, for all $i\neq j$, conditional
  independence of the leaves given the center implies that
  $\+L_i,P_i,\+L_j,P_j$ all commute.
  The two bullets follow directly.
\end{proof}

We partition $\*1^\perp$ into three orthogonal subspaces:
\begin{align*}
  \*1^\perp = K_0 \oplus K_1 \oplus K_2,
\end{align*}
where $K_0$ contains the degree-zero component $f_\emptyset$, $K_1$
contains the degree-one components $\sum_i f_{\set{i}}$, and $K_2$
contains the remaining higher-order components.
By \Cref{prop:Hoeffding-expand}, we can formally define them as
\begin{align*}
  K_0 &:= \Ker (\+L_{N(v)}), \\
  K_1 &:= \Ker (\+L_{N(v)} - \mathbb{I}), \\
  K_2 &:= \bigoplus_{k=2}^d
    \Ker (\+L_{N(v)}-k\mathbb{I}).
\end{align*}
When $d<2$, the direct sum defining $K_2$ is the zero subspace.
For convenience, let $K_+ := K_1 \oplus K_2$.
For $i\in\set{0,1,2,+}$, let $\mathbb I_i$ be the identity operator on
$K_i$.
For the rest of this section, we regard every operator as a block operator
with block indices $\set{0,1,2}$.
An index $+$ refers to the combined subspace $K_+=K_1\oplus K_2$.

\begin{definition} \label{block-matrix-via-Ki}
  For a linear map $A:\*1^\perp\to\*1^\perp$ and
  $i,j\in\set{0,1,2,+}$, define its projected block by
  \begin{align*}
    A_{ij} := \operatorname{proj}_{K_i} A \vert_{K_j}
    : K_j \longrightarrow K_i.
  \end{align*}
\end{definition}

\begin{theorem} \label{poincare-via-proj-norms}
  Suppose that $r := \norm{(P_v)_{00}} < 1/5$ and
  \[\eta := \frac{\norm{(P_v)_{10}}^2}{(1-r)(1-4r)} < 1.\]
  Then the following Poincar{\'e} inequality holds with $\eta$.
  \begin{align} \label{target:spgap-weighted-GD}
    \frac{1-\eta}{2}\inner{f}{f} \leq \inner{f}{\tp{\+L_v + \frac{1+\eta}{2}\+L_{N(v)}} f}, \quad f \in \*1^\perp.
  \end{align}
\end{theorem}

The proof of \Cref{poincare-via-proj-norms} is given in
\Cref{subsec:star-schur-complement-proof}.

\subsection{A Schur-complement proof}
\label{subsec:star-schur-complement-proof}
We now prove \Cref{poincare-via-proj-norms}.
Recall that we work in the inner-product space $(\*1^\perp, \inner{\cdot}{\cdot}_\nu)$.
Then \eqref{target:spgap-weighted-GD} can be rewritten as a matrix inequality.
It suffices to show the following
operator inequality on $(\*1^\perp,\inner{\cdot}{\cdot})$:
\begin{align*}
  G := \+L_v + \frac{1+\eta}{2}\+L_{N(v)} - \frac{1-\eta}{2} \mathbb{I} \succeq 0.
\end{align*}
By \Cref{block-matrix-via-Ki},
$G$ can be written as
\begin{align*}
  G &= \begin{bmatrix}
    (\+L_v)_{00} - \frac{1-\eta}{2}\mathbb{I}_0 & (\+L_v)_{0+} \\
    (\+L_v)_{+0} & (\+L_v)_{++} - \frac{1-\eta}{2}\mathbb{I}_+ + \frac{1+\eta}{2} (\+L_{N(v)})_{++}
  \end{bmatrix}
\end{align*}
where $\+L_{N(v)}$ vanishes on $K_0$ and preserves $K_+$ by definition.

\begin{proof}[Proof of \Cref{poincare-via-proj-norms}]
Recall that $0\leq\eta<1$.
Since $(\+L_v)_{00}=\mathbb{I}_0-(P_v)_{00}$, we have
\begin{align*}
  G_{00}=(\+L_v)_{00}-\frac{1-\eta}{2}\mathbb{I}_0
  =\frac{1+\eta}{2}\mathbb{I}_0-(P_v)_{00}\succ0.
\end{align*}
Thus \Cref{prop:Schur-complement-elimination} gives
\begin{align*}
  G \succeq 0 \quad \Longleftrightarrow \quad G/G_{00} \succeq 0.
\end{align*}
Hence it remains to prove $G/G_{00}\succeq0$.  A direct calculation gives
\begin{align*}
  G/G_{00}
  &= G_{++} - G_{+0}G_{00}^{-1}G_{0+} \\
  &= (\+L_v)_{++} - \frac{1-\eta}{2}\mathbb{I}_+
     + \frac{1+\eta}{2} (\+L_{N(v)})_{++} \\
  &\quad - (\+L_v)_{+0}
    \tp{(\+L_v)_{00} - \frac{1-\eta}{2}\mathbb{I}_0}^{-1}
    (\+L_v)_{0+} \\
  &= \+L_v/(\+L_v)_{00} - \frac{1-\eta}{2}\mathbb{I}_+ + \frac{1+\eta}{2} (\+L_{N(v)})_{++} \\
  &\quad + (\+L_v)_{+0}\left((\+L_v)_{00}^{-1}
    - \tp{(\+L_v)_{00} - \frac{1-\eta}{2}\mathbb{I}_0}^{-1}\right)
    (\+L_v)_{0+} \\
  &= \+L_v/(\+L_v)_{00} - \frac{1-\eta}{2}\mathbb{I}_+ + \frac{1+\eta}{2} (\+L_{N(v)})_{++} \\
  &\quad - \frac{1-\eta}{2} (\+L_v)_{+0}(\+L_v)_{00}^{-1}
    \tp{(\+L_v)_{00} - \frac{1-\eta}{2}\mathbb{I}_0}^{-1}
    (\+L_v)_{0+}.
\end{align*}
Since $\+L_v\succeq0$ and $(\+L_v)_{00}\succ0$, another application of
\Cref{prop:Schur-complement-elimination} gives
$\+L_v/(\+L_v)_{00}\succeq0$.  Moreover,
\begin{align*}
  (\+L_v)_{+0}(\+L_v)_{00}^{-1}
  \tp{(\+L_v)_{00} - \frac{1-\eta}{2}\mathbb{I}_0}^{-1}
  (\+L_v)_{0+}
  &\preceq
  \frac{(P_v)_{+0}(P_v)_{0+}}
  {(1-r)(\frac{1+\eta}{2}-r)}.
\end{align*}
Here $(\+L_v)_{0+}=-(P_v)_{0+}$ and
$(\+L_v)_{+0}=-(P_v)_{+0}$, so the signs cancel.  Also,
\Cref{prop:Hoeffding-expand} gives
\begin{align*}
  (\+L_{N(v)})_{++} \succeq
  \begin{bmatrix}
    \mathbb{I}_1 & 0 \\
    0 & 2 \mathbb{I}_2
  \end{bmatrix}.
\end{align*}
Combining these bounds and dropping $\eta$ from the feedback coefficient
gives
\begin{align*}
  G/G_{00} \succeq 
  \begin{bmatrix}
    \eta \mathbb{I}_1 & 0 \\
    0 & \frac{1}{2} \mathbb{I}_2
  \end{bmatrix} - \underbrace{\frac{1}{(1-r)(1-2r)}
  \begin{bmatrix}
    (P_v)_{10} \\ (P_v)_{20} 
  \end{bmatrix}
  \begin{bmatrix}
    (P_v)_{01} & (P_v)_{02} 
  \end{bmatrix} }_{=:F}.
\end{align*}
Denote the right-hand side by $N$.  We regard $F$ and $N$ as block
operators on $K_+=K_1\oplus K_2$, and
write $F_{ij}$ and $N_{ij}$ for their blocks, where $i,j\in\{1,2\}$.
We first claim that
\begin{align} \label{eq:bound-||F||}
  0 \preceq F \preceq \frac{r}{1-2r}\mathbb{I}_+.
\end{align}
Since $r<1/5$, \eqref{eq:bound-||F||} implies
$N_{22}\succeq(1/2-r/(1-2r))\mathbb{I}_2\succ0$.
  After applying the Schur complement with $K_2$ as the eliminated subspace, it
remains to show $N/N_{22}\succeq0$, where
$N/N_{22}:=N_{11}-N_{12}N_{22}^{-1}N_{21}$.  We have
\begin{align*}
  N/N_{22}
  &= \eta \mathbb{I}_1 - F_{11} - F_{12}N_{22}^{-1}F_{21} \\
  &\succeq \eta\mathbb{I}_1 - \norm{F_{11}}\mathbb{I}_1 - \frac{\norm{F_{12}}^2}{\frac{1}{2} - \norm{F}}\mathbb{I}_1.
\end{align*}
Since $F\succeq 0$, we have $\norm{F_{12}}^2 \leq \norm{F}\norm{F_{11}}$. This gives
\begin{align*}
  N/N_{22}
  &\succeq \eta\mathbb{I}_1 - \norm{F_{11}}\mathbb{I}_1 - \frac{\norm{F_{11}}\norm{F}}{\frac{1}{2} - \norm{F}}\mathbb{I}_1 \\
  &= \eta\mathbb{I}_1
  - \frac{\norm{(P_v)_{10}}^2}{(1-r)(1-2r)}
    \frac{1}{1-2\norm{F}}\mathbb{I}_1 \\ 
  &\succeq \left(\eta
    - \frac{\norm{(P_v)_{10}}^2}{(1-r)(1-4r)}\right)
    \mathbb{I}_1
  \succeq0.
\end{align*}
The last two steps use \eqref{eq:bound-||F||} and the choice of $\eta$.

It remains to prove \eqref{eq:bound-||F||}.  By definition,
\begin{align*}
  \norm{F} \leq
  \frac{\norm{(P_v)_{+0}}^2}{(1-r)(1-2r)}.
\end{align*}
Since $P_v^2=P_v$,
\begin{align*}
  (P_v)_{00}=(P_v)_{00}^2+(P_v)_{0+}(P_v)_{+0}.
\end{align*}
Since $0\preceq(P_v)_{00}\preceq\mathbb I_0$ and
$\norm{(P_v)_{00}}=r<1/2$, the spectrum of $(P_v)_{00}$ lies in $[0,r]$.
Therefore
\begin{align*}
  \norm{(P_v)_{+0}}^2
  =\norm{(P_v)_{00}-(P_v)_{00}^2}
  \leq r(1-r).
\end{align*}
Combining the last two bounds proves \eqref{eq:bound-||F||} and finishes
the proof.
\end{proof}


\section{Rapid mixing via local spectral contraction}\label{sec:general}
In this section, we will prove \Cref{thm:general}.
Let $G = (V,E)$ be a graph with maximum degree $\Delta$ and girth at least $5$.
Recall that $A \in [0,1]^{q\times q}$ is the symmetric interaction matrix and $\*\lambda \in \mathbb{R}_{\geq 0}^q$ is the external field.
The Gibbs distribution $\mu_G$ with interaction matrix $A$ and external field $\*\lambda$ on $G$ is defined by
\begin{align*}
  \mu_G(\sigma) \propto \prod_{uv \in E} A_{\sigma_u,\sigma_v} \prod_{v\in V} \lambda_{\sigma_v}, \quad \forall \sigma \in [q]^V.
\end{align*}

Fix $v \in V$ and a feasible pinning $\tau \in [q]^{V\setminus N[v]}$.
Then, for convenience, let
\begin{align} \label{eq:def-nu-pi}
  \nu &:= \mu^\tau_{G,N[v]},
  \quad
  \nu_v := \mu^\tau_{G,v},
  \quad \text{and} \quad
  p_i := \mu^\tau_{G-v, i}, \quad \forall i\in N(v).
\end{align}
We also identify $N(v)$ with $[d]$, where $d = \abs{N(v)}$.

It follows directly that \Cref{cond-local-spectral-contraction} implies the following condition on $\nu, (p_i)_{i=1}^d$.

\begin{condition}\label{cond:alpha-eps-Delta-contraction}
  There are $\alpha \geq 1$ and $\epsilon \in (0,1]$ such that, with
  %Let $\alpha \geq 1$, $\epsilon \in (0,1]$, and $\Delta \geq d \geq 1$.
  %We say $(A,\*\lambda, (\*p_i)_{i=1}^d)$ satisfies \emph{$(\alpha, \epsilon, \Delta)$-contraction} if
  %$A$ is symmetric and $0 \leq A \leq 1$ entry-wise such that, let
  \[B := \*1\*1^\intercal - A,\]
  the following conditions hold for $\nu_v$ and $(p_i)_{i=1}^d$:
  \begin{itemize}
  \item ($\alpha$-marginal bound) for every $i\in[d]$ and
    $\!c\in[q]$,
    \begin{align*}
      (Bp_i)(\!c)\leq \frac{\alpha}{\Delta}, \quad (B\nu_v)(\!c)\leq \frac{\alpha}{\Delta};
    \end{align*}
  \item ($\epsilon$-contraction) we have
    \begin{align*}
      \norm{
      \-{diag}\tp{\sum_{i=1}^d p_i}^{1/2}
      B
      \-{diag}\tp{\nu_v}^{1/2}
      }_{2\to2}^2 \leq \frac{1-\epsilon}{\Delta}.
    \end{align*}
  \end{itemize}
\end{condition}

The main result of this section is the following lemma.

\begin{lemma}\label{lem:block-estimate-via-contraction}
  Suppose \Cref{cond:alpha-eps-Delta-contraction} holds and $\Delta\geq2\alpha$.
  Then, on the inner-product space $(\*1^\perp, \inner{\cdot}{\cdot}_{\nu})$, we have
  \begin{align}
    \label{eq:general-P00}
    \norm{(P_v)_{00}} &\leq O_{\alpha}\tp{1/\Delta}, \\
    \label{eq:general-P10}
    \norm{(P_v)_{10}}^2 &\leq \tp{1 + O_\alpha(1/\Delta)} \cdot \frac{1-\epsilon}{\Delta} + O_\alpha\tp{\frac{\sqrt{\log(2d)}}{\Delta^{3/2}}}.
  \end{align}
\end{lemma}

We are now ready to prove \Cref{thm:general}.
\begin{proof}[Proof of \Cref{thm:general}]
For every $v\in V$, set
\begin{align*}
  \eta_v:=\frac{1-\epsilon/2}{\Delta}.
\end{align*}
We verify the two hypotheses of \Cref{thm:local-to-global} for these
parameters.

Fix $v\in V$ and a feasible pinning $\tau$ outside $N[v]$, and let
$d=\abs{N(v)}$.  If $d=0$, then the weighted dynamics on $N[v]=\set{v}$
has Laplacian $\+L_v$ and satisfies the Poincar{\'e} inequality with
constant one, which is at least $(1-\eta_v)/2$.  We may therefore assume
that $d\geq1$.

Use the notation in \eqref{eq:def-nu-pi}.  Since $G$ has girth at least
$5$, the induced graph on $N[v]$ is a star, and, conditional on the
center, its leaves are independent.  Moreover,
\Cref{cond-local-spectral-contraction} implies
\Cref{cond:alpha-eps-Delta-contraction}.  Hence
\Cref{lem:block-estimate-via-contraction} gives
\begin{align*}
  r:=\norm{(P_v)_{00}}
  &=O_\alpha\tp{\frac1\Delta},
  \\
  \norm{(P_v)_{10}}^2
  &\leq
  \tp{1+O_\alpha\tp{\frac1\Delta}}
  \frac{1-\epsilon}{\Delta}
  +O_\alpha\tp{
    \frac{\sqrt{\log(2d)}}{\Delta^{3/2}}
  }.
\end{align*}
In particular, by taking the lower bound on $\Delta$ sufficiently large
in terms of $\alpha$ and $\epsilon$, we have $r<1/5$ and
\begin{align*}
  \widehat\eta
  &:=\frac{\norm{(P_v)_{10}}^2}{(1-r)(1-4r)}
  \\
  &\leq
  \frac{1-\epsilon}{\Delta}
  +O_\alpha\tp{
    \frac1{\Delta^2}
    +\frac{\sqrt{\log(2\Delta)}}{\Delta^{3/2}}
  }
  \leq
  \frac{1-\epsilon/2}{\Delta}
  =\eta_v<1,
\end{align*}
where we used $d\leq\Delta$ and $r=O_\alpha(1/\Delta)$ in the first
inequality.

Applying \Cref{poincare-via-proj-norms} with $\widehat\eta$ shows that,
for every $f\perp\*1$ in $L^2(\nu)$,
\begin{align*}
  \frac{1-\widehat\eta}{2}\norm{f}^2
  \leq
  \inner{f}{\tp{
    \+L_v+\frac{1+\widehat\eta}{2}\+L_{N(v)}
  }f}.
\end{align*}
Since $\widehat\eta\leq\eta_v$ and $\+L_{N(v)}\succeq0$, this implies
\begin{align*}
  \frac{1-\eta_v}{2}\norm{f}^2
  \leq
  \inner{f}{\tp{
    \+L_v+\frac{1+\eta_v}{2}\+L_{N(v)}
  }f}.
\end{align*}
Thus the required weighted spectral-gap estimate holds uniformly over
all feasible pinnings $\tau$.

For every $v\in V$, the maximum-degree assumption gives
\begin{align*}
  \sum_{u\in N(v)}\eta_u
  \leq
  \Delta\cdot\frac{1-\epsilon/2}{\Delta}
  =1-\frac\epsilon2.
\end{align*}
The factorization defining $\mu_G$ also implies the conditional
independence property in \Cref{def:cond-indep}.  Therefore
\Cref{thm:local-to-global}, applied with global reserve $\epsilon/2$,
shows that the continuous-time Glauber Laplacian $\+L$ has spectral gap
at least $\epsilon/4$.

Finally, the Laplacian of the discrete-time Glauber dynamics is
$\+L/n$.  Its spectral gap is consequently at least
$\epsilon/(4n)$, which proves the theorem.
\end{proof}

In the remainder of this section, we prove \Cref{lem:block-estimate-via-contraction}.
In particular, \eqref{eq:general-P00} and \eqref{eq:general-P10} will be proved in \Cref{sec:general-P00} and \Cref{sec:general-P10}, respectively.

We will work in the inner product space
\begin{align*}
  (\*1^\perp, \inner{\cdot}{\cdot}_\nu).
\end{align*}
We will first introduce some basic definitions and facts.
By the definition of $\nu$, we have
%
%Let $\*a \in \mathbb{R}^q_{\geq 0}$ be a vector.
%Let $A \in \mathbb{R}^{q\times q}_{\geq 0}$ be a symmetry matrix.
%\begin{align*}
%  A := \-{diag}\set{\*a} + \*1\*1^\intercal.
%\end{align*}
%Let $\*\lambda \in \mathbb{R}^q_{\geq 0}$.
%For every $i$, let $p_i \in \mathbb{R}^q_{\geq 0}$ be a distribution on $[q]$.
%We use
%\[\rho := \otimes_i p_i,\]
%to denote the product distribution on leaves.
%We define the joint distribution $\nu$ on star:
\begin{align} \label{eq:def-local-mu}
  \nu(v\!c, \sigma) \propto \lambda_{\!c}\prod_{i}A_{\!c,\sigma_i}p_i(\sigma_i), \quad \forall \!c \in [q], \sigma \in [q]^d.
\end{align}
In particular, this implies the \emph{tree recursion}:
\begin{align*}
  \nu(v\!c) = g_{\!c}(\*p) := \frac{\lambda_{\!c} \prod_i (Ap_i)(\!c)}{\sum_{\!b} \lambda_{\!b} \prod_i (Ap_i)(\!b)}.
\end{align*}
If the center $v$ is fixed to color $\!c$, then the posterior distribution on leaf $i$ is
\begin{align*}
  \nu^{v\!c}(i\!a) := \frac{A_{\!c,\!a}p_i(\!a)}{(Ap_i)(\!c)}, \quad \forall \!a \in [q].
\end{align*}

\begin{definition}\label{def:L}
  For the product distribution $\rho = \otimes_i p_i$ on the leaves, define the following function on $\-{supp}(\rho)$:
  \begin{align*}
    L_\rho(\tau) := \frac{\nu(\tau)}{\rho(\tau)}, \quad \forall \tau \in \-{supp}(\rho).
  \end{align*}
  When the context is clear, we drop the subscript $\rho$ from $L_\rho$ and directly use $L$.

  Moreover, we define similar ratios for conditional distributions:
  \begin{align*}
    \forall \!c,i, \quad L^{\!c}_i(\!a) &:= \frac{\nu^{v\!c}(i\!a)}{\rho_i(\!a)} = \frac{A_{\!c,\!a}}{(Ap_i)_{\!c}}, \quad \!a\in [q],\\
    \forall \!c, \quad L^{\!c}(\tau) &:= \prod_i L_i^{\!c}(\tau_i) = \prod_{j=1}^d \frac{A_{\!c,\tau_j}}{(Ap_j)(\!c)}, \quad \forall \tau \in \mathrm{supp}(\rho).
  \end{align*}
  For a fixed $i$, we measure the influence from $v$ to $i$ by $\Psi_i \in \mathbb{R}^{q\times q}$:
\begin{align*}
  \Psi_{i}(\!c,\!d) := \mathrm{Cov}_{p_i}(L^{\!c}_i,L^{\!d}_i).
\end{align*}
\end{definition}

% \begin{remark}
%   The function $L$ is useful in the sense that we can use it to approximate the distribution $\nu$ on leaves by the product distribution $\rho$.
%   For example, let $(f_i)_{i=1}^d$ be a family of functions in $\mathbb{R}^q$. Then,
%   \begin{align*}
%     \E[\nu]{\prod_i f_i}
%     &= \E[\nu]{(\prod_i f_i) L } + \E[\nu]{(\prod_i f_i) (1 - L)} \approx \prod_i \E[\rho_i]{f_i},
%     %&= \E[\rho]{\prod_i f_i} + \E[\!c\sim\nu_v]{\E[\nu^{v\!c}]{(\prod_i f_i) (1 - L)}} \approx \prod_i \E[\rho_i]{f_i},
%   \end{align*}
%   where the last equation is justified by the following concentration bound on $L$.
% \end{remark}

\begin{proposition} \label{prop:Psi-i}
  For each fixed $i$,
  \begin{align*}
    \Psi_i = \mathrm{diag}(\*1-Bp_i)^{-1}B\tp{\mathrm{diag}(p_i) - p_ip_i^\intercal}B \, \mathrm{diag}(\*1-Bp_i)^{-1}.
  \end{align*}
\end{proposition}
\begin{proof}
  Fix $\!c,\!d$. By definition, we have
  \begin{align*}
    \Psi_i(\!c,\!d)
    &= \E[p_i]{L^{\!c}_i L^{\!d}_i} - 1 
    =   \frac{\sum_{\!b}p_i(\!b)A_{\!c,\!b}A_{\!d,\!b}}{(Ap_i)(\!c) (Ap_i)(\!d)} - 1 \\
    \tag{$B = \*1\*1^\intercal - A$}
    &=   \frac{\sum_{\!b}p_i(\!b)B_{\!c,\!b}B_{\!d,\!b} -(Bp_i)(\!c)(Bp_i)(\!d)}{(1-(Bp_i)(\!c))(1 - (Bp_i)(\!d))}.
  \end{align*}
  This finishes the proof.
\end{proof}

\begin{corollary} \label{cor:Psi-related-bounds}
  Under \Cref{cond:alpha-eps-Delta-contraction}, if $\Delta\geq2\alpha$, then
  \begin{itemize}
  \item for every $i$, $\!c,\!d$, $\abs{\Psi_i(\!c,\!d)} = O_\alpha(\Delta^{-1})$;
  \item for every $i$, $\norm{\abs{\Psi_i} \nu_v}_\infty = O_{\alpha}(\Delta^{-2})$,
  \end{itemize}
  where $\norm{\cdot}_\infty$ is the Euclidean norm, and $\abs{\cdot}$ means taking entrywise absolute values.
\end{corollary}
\begin{proof}
  For the first bound, by the $\alpha$-marginal bound, \Cref{prop:Psi-i}, and $B \leq 1$ entry-wise, %we have
  \begin{align*}
    \abs{\Psi_i(\!c,\!d)}
    &\leq (1-\alpha/\Delta)^{-2}\tp{(Bp_i)(\!c) + \frac{\alpha^2}{\Delta^2}} \\
    &\leq (1-\alpha/\Delta)^{-2}\tp{\frac{\alpha}{\Delta} + \frac{\alpha^2}{\Delta^2}} = O_\alpha(\Delta^{-1}).
  \end{align*}
  Similarly, for the second bound, fix $\!c$. By the $\alpha$-marginal bound and \Cref{prop:Psi-i}, we have
  \begin{align*}
    (\abs{\Psi_i} \nu_v)(\!c)
    &= \sum_{\!d} \nu(v\!d)\abs{\Psi_i(\!c,\!d)} \\
    &\leq (1-\alpha/\Delta)^{-2}\sum_{\!d} \nu(v\!d)\tp{\sum_{\!b}p_i(\!b)B_{\!c,\!b}B_{\!d,\!b} + \frac{\alpha^2}{\Delta^2}} \\
    &=(1-\alpha/\Delta)^{-2} \sum_{\!b}p_i(\!b)B_{\!c,\!b} (B\nu_v)(\!b) + O_\alpha(\Delta^{-2}) \\
    &\leq (1-\alpha/\Delta)^{-2} \cdot \frac{\alpha}{\Delta} (Bp_i)(\!c) + O_\alpha(\Delta^{-2}) = O_\alpha(\Delta^{-2}). \qedhere
    %&\leq (1-\alpha/\Delta)^{-2} \cdot \frac{\alpha^2}{\Delta^2} + O_\alpha(\Delta^{-2}) = O_\alpha(\Delta^{-2}). \qedhere
  \end{align*}
\end{proof}

\begin{fact}\label{fact:var-matrix-form}
  Let $\pi$ be a distribution on $[q]$.
  Let $f:[q] \to \mathbb{R}$ be a function.
  Then
  \begin{align*}
    \Var[\pi]{f} = f^\intercal \tp{\mathrm{diag}(\pi) - \pi\pi^\intercal} f,
  \end{align*}
  where we treat $\pi$ and $f$ as vectors in $\mathbb{R}^q$.
\end{fact}

\begin{proposition}\label{prop:normalizer-conditional-tail}
Suppose $1\leq d\leq\Delta$ and the $\alpha$-marginal bound in
\Cref{cond:alpha-eps-Delta-contraction} holds.  Let
$\rho=\bigotimes_{i=1}^d p_i$ and $L=L_\rho$.
If $\Delta\geq2\alpha$, then there
are constants $C_{\alpha},c_{\alpha}>0$ such that, for
every $t\geq2C_{\alpha}d/\Delta^2$,
\begin{align}
  \sup_{\!r:\,\nu(v\!r)>0}
  \Pr[\nu^{v\!r}]{\abs{L-1}>t}
  \leq
  2\exp\tp{-c_{\alpha}\frac{\Delta^2}{d}t^2}
  \leq
  2\exp\tp{-c_{\alpha}\Delta t^2}.
  \label{eq:normalizer-conditional-tail}
\end{align}
\end{proposition}

\begin{proof}[Proof of \Cref{prop:normalizer-conditional-tail}]
%Let $B:=\*1\*1^\intercal-A$ as in
%\Cref{cond:alpha-eps-Delta-contraction}, and set
%\begin{align*}
%  \pi_{\!c}:=\nu(v\!c),
%  \qquad
%  z_{i,\!c}:=A_{\!c}p_i,
%  \qquad
%  \eta:=\frac{\alpha}{\Delta}.
%\end{align*}
Center colors of zero $\nu$-mass are omitted below.
%
For every $i$ and $\!c$, by the $\alpha$-marginal bound and $\Delta \geq 2\alpha$,
\begin{align}
  1-\frac{\alpha}{\Delta}
  \leq
  (Ap_i)(\!c)
  %z_{i,\!c}
  =1-(Bp_i)(\!c)
  \leq 1,
  \qquad
  \frac{\alpha}{\Delta}\leq\frac12.
  \label{eq:normalizer-z-bounds}
\end{align}
%For $x\in\-{supp}(\rho)$, define
%\begin{align*}
%  \ell_{\!c}(x)
%  :=
%  \prod_{i=1}^d\frac{A_{\!c,x_i}}{(Ap_i)(\!c)}.
%\end{align*}
%By \eqref{eq:def-local-mu}, $\ell_{\!c}$ is the density of
%$\nu^{v\!c}$ with respect to $\rho$, and
%\begin{align}
%  L=\sum_{\!c}\nu(v\!c)\ell_{\!c}.
%  \label{eq:normalizer-likelihood-mixture}
%\end{align}
Moreover, \eqref{eq:normalizer-z-bounds} and
$0\leq A_{\!c,\!a}\leq1$ give
\begin{align}
  0\leq L^{\!c},L
  \leq
  \tp{1-\frac{\alpha}{\Delta}}^{-d}
  \leq
  \exp\tp{\frac{2\alpha d}{\Delta}}
  \leq
  \exp(2\alpha).
  \label{eq:normalizer-uniform-bound}
\end{align}

Suppose that two leaf configurations differ only at coordinate $i$,
where the leaf color changes from $\!a$ to $\!b$.  Then
\begin{align*}
  \abs{L^{\!c}(x)-L^{\!c}(x')}
  &\leq
  \tp{1-\frac{\alpha}{\Delta}}^{-d}
  \abs{A_{\!c,\!a}-A_{\!c,\!b}}
  \\
  &\leq
  \tp{1-\frac{\alpha}{\Delta}}^{-d}
  \tp{B_{\!c,\!a}+B_{\!c,\!b}}.
\end{align*}
Thus the symmetry of $B$ and the center interaction bound imply
\begin{align}
  \abs{L(x)-L(x')}
  &\leq
  \tp{1-\frac{\alpha}{\Delta}}^{-d}
  \sum_{\!c}\nu(v\!c)
  \tp{B_{\!c,\!a}+B_{\!c,\!b}}
  \notag\\
  \notag
  &= \tp{1-\frac{\alpha}{\Delta}}^{-d}\tp{(B\nu_v)(\!a) + (B\nu_v)(\!b)}\\
  &\leq
  \frac{2\alpha\exp(2\alpha)}{\Delta}.
  \label{eq:normalizer-bounded-difference}
\end{align}
We claim that there is $C_\alpha > 0$ such that for every $\!d$,
\begin{align}
  \abs{\E[\nu^{v\!d}]{L}-1}
  \leq
  \frac{C_{\alpha}d}{\Delta^2}.
  \label{eq:normalizer-conditional-mean-target}
\end{align}
Since the leaves are independent under every $\nu^{v\!d}$, \Cref{prop:normalizer-conditional-tail} now follows from \eqref{eq:normalizer-bounded-difference},
\eqref{eq:normalizer-conditional-mean-target}, and McDiarmid's inequality;
the second inequality in \eqref{eq:normalizer-conditional-tail} uses
$d\leq\Delta$.

It only remains to prove \eqref{eq:normalizer-conditional-mean-target}.
Note that 
\begin{align*}
  \E[\nu^{v\!d}]{L}
  &= \sum_{x} \nu^{v\!d}(x) \sum_{\!c} \nu(v\!c) L^{\!c}(x) \\
  &= \sum_{\!c}\nu(v\!c) \sum_{x} \rho(x) L^{\!d}(x) L^{\!c}(x) \\
  \tag{independent}
  &= \sum_{\!c}\nu(v\!c) \prod_{i} (1 + \Psi_i(\!d,\!c)).
\end{align*}
This implies that
\begin{align*}
  \abs{\E[\nu^{v\!d}]{L}-1}
  %\nonumber
  %&= \abs{\sum_{\!c}\nu(v\!c)(\prod_{i} (1 + \Psi_i(\!d,\!c))-1)} \\
  \tag{triangle ineq.}
  &\leq \sum_{\!c}\nu(v\!c)\abs{\prod_{i} (1 + \Psi_i(\!d,\!c))-1)} \\
  %\label{eq:diff-expectation}
  &\leq \sum_{\!c}\nu(v\!c)\exp\tp{\sum_i \abs{\Psi_i(\!d,\!c)}} \sum_i \abs{\Psi_i(\!d,\!c)} \\
  \tag{\Cref{cor:Psi-related-bounds}, 1st bullet}
  &= \sum_i(\abs{\Psi_i}\nu_v)(\!d) \cdot O_\alpha(1) \\
  \tag{\Cref{cor:Psi-related-bounds}, 2nd bullet}
  &= \frac{O_\alpha(1)d}{\Delta^2}.
\end{align*}
Here the second inequality uses the following bound:
\begin{align*}
  \abs{\prod_i(1+x_i)-1}
  \leq
  \exp\tp{\sum_i\abs{x_i}}
  \sum_i\abs{x_i}.
\end{align*}
This finishes the proof.
%By \Cref{prop:Psi-i}, we have
%\begin{align*}
%  \abs{\Psi_i(\!d,\!c)}
%  &= \frac{\sum_{\!b}p_j(\!b)B_{\!c,\!b}B_{\!d,\!b} -(Bp_j)(\!c)(Bp_j)(\!d)}{(1-(Bp_j)(\!c))(1 - (Bp_j)(\!d))} \\
%  \tag{by \eqref{eq:normalizer-z-bounds}}
%  &\leq 
%    4\tp{
%    \sum_{\!b}p_i(\!a)B_{\!c,\!b}B_{\!d,\!b}
%    +\frac{\alpha^2}{\Delta^2}
%    }. 
%\end{align*}
%Finally, summing over $i$, we have
%\begin{align*}
%  \sum_i \abs{\Psi_i(\!d,\!c)} &\leq 4\sum_i
%\end{align*}
%
%\begin{align}
%  \abs{\E[\nu^{v\!r}]{L}-1}
%  &=
%  \abs{\sum_{\!c}\nu(v\!c)(K_{\!r,\!c}-1)}
%  \leq
%  \frac{C_{\alpha}d}{\Delta^2}.
%  \label{eq:normalizer-conditional-mean}
%\end{align}
%
%
%For center colors $\!c,\!r$, set
%\begin{align*}
%  K_{\!c,\!r}
%  &:=
%  \E[\rho]{\ell_{\!c}\ell_{\!r}},
%  &
%  W_{\!c,\!r}
%  &:=
%  \sum_{i=1}^d\sum_{\!a}
%  p_i(\!a)B_{\!c,\!a}B_{\!r,\!a}.
%\end{align*}
%Independence under $\rho$ gives
%\begin{align}
%  K_{\!c,\!r}
%  =
%  \prod_{i=1}^d\tp{1+\gamma_i(\!c,\!r)},
%  \label{eq:normalizer-two-replica-product}
%\end{align}
%where
%\begin{align*}
%  \gamma_i(\!c,\!r)
%  :=
%  \frac{
%    \sum_{\!a}p_i(\!a)B_{\!c,\!a}B_{\!r,\!a}
%    -(Bp_i)(\!c)(Bp_i)(\!r)
%  }{
%    z_{i,\!c}z_{i,\!r}
%  }.
%\end{align*}
%By \eqref{eq:normalizer-z-bounds},
%\begin{align*}
%  \abs{\gamma_i(\!c,\!r)}
%  \leq
%  4\tp{
%    \sum_{\!a}p_i(\!a)B_{\!c,\!a}B_{\!r,\!a}
%    +\frac{\alpha^2}{\Delta^2}
%  }.
%\end{align*}
%Also,
%\begin{align*}
%  0\leq W_{\!c,\!r}
%  \leq
%  \sum_{i=1}^d(Bp_i)(\!r)
%  \leq\frac{\alpha d}{\Delta}
%  \leq\alpha.
%\end{align*}
%The elementary product bound
%\begin{align*}
%  \abs{\prod_i(1+x_i)-1}
%  \leq
%  \exp\tp{\sum_i\abs{x_i}}
%  \sum_i\abs{x_i}
%\end{align*}
%applied to \eqref{eq:normalizer-two-replica-product} yields
%\begin{align}
%  \abs{K_{\!c,\!r}-1}
%  \leq
%  C_{\alpha}
%  \tp{W_{\!c,\!r}+\frac{d}{\Delta^2}}.
%  \label{eq:normalizer-overlap-bound}
%\end{align}
%Furthermore, symmetry and the two interaction bounds give
%\begin{align}
%  \sum_{\!c}\pi_{\!c}W_{\!r,\!c}
%  &=
%  \sum_{i=1}^d\sum_{\!a}
%  p_i(\!a)B_{\!r,\!a}(B\nu_v)(\!a)
%  \notag\\
%  &\leq
%  \frac{\alpha}{\Delta}
%  \sum_{i=1}^d(Bp_i)(\!r)
%  \leq
%  \frac{\alpha^2d}{\Delta^2}.
%  \label{eq:normalizer-W-sum}
%\end{align}
%Consequently,
%\begin{align}
%  \abs{\E[\nu^{v\!r}]{L}-1}
%  &=
%  \abs{\sum_{\!c}\pi_{\!c}(K_{\!r,\!c}-1)}
%  \leq
%  \frac{C_{\alpha}d}{\Delta^2}.
%  \label{eq:normalizer-conditional-mean}
%\end{align}
%
\end{proof}


%In the rest part of this section, we will work in the space $L^2(\nu)$.


\subsection{Bound for $\norm{(P_v)_{00}}$} \label{sec:general-P00}
In this subsection, we prove \eqref{eq:general-P00}.
We will need the following result.
\begin{observation} \label{obs:PvhL-on-i}
  For $h \in K_0$, let $F := P_v h \cdot L$ be a function on the leaves.
  Fix $i\in[d]$ and $x \in \mathrm{supp}(\rho)$. Then
  \begin{align*}
    F(x_{\sim i}, x_i) = (Aw_{i,x_{\sim i}})(x_i),
  \end{align*}
  where the vector $w_{i,x_{\sim i}} \in \mathbb{R}^q$ is defined as
  \begin{align*}
    w_{i,x_{\sim i}}(\!c)
    :=
    \frac{\nu_v(\!c)h(\!c)}{(Ap_i)(\!c)}
    \prod_{j\ne i}L^{\!c}_j,\quad \forall \!c.
  \end{align*}
\end{observation}
\begin{proof}
  We have
  \begin{align*}
    F(x)
    := P_vh(x)L(x) &=\sum_{\!c}\nu_v(\!c)h(\!c)L^{\!c}(x) \\
    &= \sum_{\!c} A_{x_i,\!c} \cdot \frac{\nu(v\!c)h(\!c)}{(Ap_i)(\!c)} \prod_{j\neq i} L^{\!c}_j. \qedhere
  \end{align*}
\end{proof}

\begin{proof}[Proof of \eqref{eq:general-P00}]
  %\todo{notation $\pi_{\!c}$, $z_{i,\!c}$ are unnecessary for this proof}
If $d=0$, then $P_vh=\E[\nu]{h}=0$ for every $h\in K_0$, and the
claim is immediate.  We henceforth assume that $d\geq1$.
Since $\norm{(P_v)_{00}}\leq1$, values of $\Delta$ bounded in terms of
$\alpha$ can be absorbed into the implicit constant.  We may therefore
assume throughout that $\Delta$ is sufficiently large in terms of
$\alpha$; in particular, \Cref{prop:normalizer-conditional-tail}
applies with $t=1/2$.

Because $P_v$ is an orthogonal projection,
\begin{align*}
  \inner{h}{(P_v)_{00}h}
  =\inner{h}{P_vh}
  =\norm{P_vh}^2.
\end{align*}
Split this last norm according to the events $L\geq1/2$ and $L<1/2$.
On the first event, we claim the following bound
\begin{align} \label{eq:general-P00-second-moment-target}
  \E[\rho]{(P_vh \cdot L)^2} \leq C_\alpha\frac{d}{\Delta^2} \norm{h}^2.
\end{align}
%changing measure from the leaf marginal of $\nu$
%to $\rho$ and using \eqref{eq:general-P00-second-moment} give
This implies that
\begin{align*} 
  \E[\nu]{
    \*1_{\set{L\geq1/2}}(P_vh)^2
  }
  =
  \E[\rho]{
    \*1_{\set{L\geq1/2}}\frac{(P_vh \cdot L)^2}{L}
  }
  \leq
  C_\alpha\frac{d}{\Delta^2}\norm{h}^2.
\end{align*}
On the complementary event, conditional Jensen gives
$(P_vh)^2\leq P_v(h^2)$.  Therefore,
\Cref{prop:normalizer-conditional-tail}, applied with $t=1/2$, gives
\begin{align*}
  \E[\nu]{
    \*1_{\set{L<1/2}}(P_vh)^2
  }
  &\leq
  \sum_{\!c}\nu(v\!c) h(\!c)^2
  \Pr[\nu^{v\!c}]{L<1/2}
  \\
  &\leq
  2\exp\tp{-c_\alpha\frac{\Delta^2}{d}}
  \norm{h}^2.
\end{align*}
Combining the two events and taking the supremum over nonzero
$h\in K_0$, we have
\begin{align*}
  \norm{(P_v)_{00}}
  \leq
  C_\alpha\frac{d}{\Delta^2}
  +2\exp\tp{-c_\alpha\frac{\Delta^2}{d}}
  =O_\alpha\tp{\frac1\Delta},
\end{align*}
where the last equality uses $d\leq\Delta$.  This proves
\eqref{eq:general-P00}.


It only remains to prove \eqref{eq:general-P00-second-moment-target}.
%
Fix $h\in K_0$.  Since $h$ depends only on the center, we write
$h(\!c)$ for its value at the center color $\!c$.  Set
\begin{align*}
  \rho&:=\bigotimes_{i=1}^d p_i,
  %&
  %\pi_{\!c}&:=\nu(v\!c),
  %&
  %z_{i,\!c}&:=A_{\!c}p_i.
\end{align*}
Center colors of zero $\nu$-mass are omitted below.  The
$\alpha$-marginal bound in
\Cref{cond:alpha-eps-Delta-contraction} gives
\begin{align}
  1-\frac{\alpha}{\Delta}
  \leq
  (Ap_i)(\!c)
  =1-(Bp_i)(\!c)
  \leq1.
  \label{eq:general-P00-z-bound}
\end{align}
For $x\in\-{supp}(\rho)$, recall
\begin{align*}
  %L^{\!c}(x)
  %&:={}
  %\prod_{j=1}^d\frac{A_{\!c,x_j}}{(Ap_j)(\!c)},
  %&
  F(x)
  &:={}
  P_vh(x)L(x)
  =\sum_{\!c}\nu_v(\!c)h(\!c)L^{\!c}(x).
\end{align*}
%The last identity follows from \eqref{eq:def-local-mu},
%\eqref{def-Pi}, and \Cref{def:L}.
Since $h$ is centered,
\begin{align*}
  \E[\rho]{F}
  =\sum_{\!c}\nu_v(\!c)h(\!c)
  =0.
\end{align*}
Tensorization of variance under $\rho$ therefore yields
\begin{align}
  \E[\rho]{F^2}
  \leq
  \sum_{i=1}^d
  \E[\rho]{
    \Var[\rho]{F\mid X_{\sim i}}
  }.
  \label{eq:general-P00-tensorization}
\end{align}

By \Cref{obs:PvhL-on-i}, we have
%Fix $i\in[d]$ and a configuration $x_{\sim i}$ of the other leaves,
%and define
%\begin{align*}
%  w_{i,x_{\sim i}}(\!c)
%  :=
%  \frac{\nu_v(\!c)h(\!c)}{(Ap_i)(\!c)}
%  \prod_{j\ne i}\frac{A_{\!c,x_j}}{(Ap_j)(\!c)}.
%\end{align*}
%As a function of $X_i$, we have
\begin{align*}
  F(x_{\sim i}, x_i) =(Aw_{i,x_{\sim i}})(x_i), \quad \text{where} \quad
  w_{i,x_{\sim i}}(\!c)
  :=
  \frac{\nu_v(\!c)h(\!c)}{(Ap_i)(\!c)}
  \prod_{j\ne i}L^{\!c}_j,\quad \forall \!c.
\end{align*}
%where we used the symmetry of $A$.
Applying
\Cref{fact:var-matrix-form} with distribution $p_i$ and function
$Aw_{i,x_{\sim i}}$, we obtain
\begin{align}
  \nonumber
  &\Var[\rho]{F\mid X_{\sim i}=x_{\sim i}}
  \\
  \nonumber
  &\qquad=
  w_{i,x_{\sim i}}^\intercal
  A\tp{\-{diag}(p_i)-p_ip_i^\intercal}A
  w_{i,x_{\sim i}}
  \\
  \nonumber
  &\qquad=
  w_{i,x_{\sim i}}^\intercal
  B\tp{\-{diag}(p_i)-p_ip_i^\intercal}B
  w_{i,x_{\sim i}}
  \\
  &\qquad\leq
  w_{i,x_{\sim i}}^\intercal
  B\-{diag}(p_i)B
  w_{i,x_{\sim i}}.
  \label{eq:general-P00-tensorization-B}  
\end{align}
Indeed, the middle identity uses
$A=\*1\*1^\intercal-B$ and
$\tp{\-{diag}(p_i)-p_ip_i^\intercal}\*1=0$.

Note that \eqref{eq:general-P00-z-bound} and $0\leq A\leq1$ give
\begin{align}
  \abs{L^{\!c}(x)}
  \leq
  \tp{1-\frac{\alpha}{\Delta}}^{-d}
  \leq\exp(2\alpha).
  \label{eq:general-P00-likelihood-bound}
\end{align}
%
Let $y(\!c):=\nu_v(\!c)\abs{h(\!c)}$.  By
\eqref{eq:general-P00-likelihood-bound},
\begin{align*}
  \abs{w_{i,x_{\sim i}}(\!c)}
  \leq\exp(2\alpha)y(\!c).
\end{align*}
Since $B\-{diag}(p_i)B$ has nonnegative entries,
\eqref{eq:general-P00-tensorization-B} and \eqref{eq:general-P00-tensorization} give
\begin{align}
  \E[\rho]{F^2}
  \nonumber
  &\leq
  \exp(4\alpha)
  y^\intercal
  B\-{diag}\tp{\sum_{i=1}^d p_i}B
  y \\
  \nonumber
  &=
  \exp(4\alpha)
  \abs{h}^\intercal\mathrm{diag}(\nu_v)
  B\-{diag}\tp{\sum_{i=1}^d p_i}B
  \mathrm{diag}(\nu_v)\abs{h} \\
  &\leq 
    \exp(4\alpha) \cdot
    \norm{
    \-{diag}(s)^{1/2}
    B
    \-{diag}(\nu_v)^{1/2}
    }_{2}^2 \norm{h}^2,
  \label{eq:general-P00-second-moment-preliminary}
\end{align}
where the $2$-norm is the Euclidean norm.

Set $s:=\sum_{i=1}^d p_i$.  The $\alpha$-marginal bound implies
\begin{align*}
  (Bs)(\!c)
  \leq\frac{\alpha d}{\Delta},
  \qquad
  (B\nu_v)(\!c)
  \leq\frac{\alpha}{\Delta},
  \quad \forall \!c \in [q].
\end{align*}
%for every $\!c$.
Thus, for every $x\in\mathbb R^q$, weighted
Cauchy--Schwarz and the symmetry of $B$ give
\begin{align*}
  &\norm{
    \-{diag}(s)^{1/2}
    B
    \-{diag}(\nu_v)^{1/2}x
  }_2^2
  \\
  &\qquad=
  \sum_{\!a}s(\!a)
  \tp{
    \sum_{\!c}B_{\!a,\!c}
    \sqrt{\nu_v(\!c)}x(\!c)
  }^2
  \\
  &\qquad\leq
  \sum_{\!a}s(\!a)(B\nu_v)(\!a)
  \sum_{\!c}B_{\!a,\!c}x(\!c)^2
  \\
  &\qquad\leq
  \frac{\alpha}{\Delta}
  \sum_{\!c}(Bs)(\!c)x(\!c)^2
  \leq
  \frac{\alpha^2d}{\Delta^2}\norm{x}_2^2.
\end{align*}
%The vector with coordinates
%$\sqrt{\nu_v(\!c)}\abs{h(\!c)}$ has squared Euclidean norm
%$\norm{h}^2$.
Hence
\eqref{eq:general-P00-second-moment-preliminary} yields
\begin{align}
  \E[\rho]{F^2}
  \leq
  C_\alpha\frac{d}{\Delta^2}\norm{h}^2,
  \label{eq:general-P00-second-moment}
\end{align}
as claimed in \eqref{eq:general-P00-second-moment-target}.
\end{proof}



\subsection{Bound for $\norm{(P_v)_{10}}$} \label{sec:general-P10}
In this subsection, we prove \eqref{eq:general-P10}.
Let $\rho=\bigotimes_{i=1}^d p_i$, and let $L=L_\rho$.
For $h \in K_0$, we have
\begin{align*}
  \norm{(P_v)_{10}h}
  &= \norm{\-{proj}_{K_1} P_v h} \\
  \tag{triangle ineq.}
  &\leq \norm{\-{proj}_{K_1} (P_v h \cdot L)} + \norm{\-{proj}_{K_1} (P_v h \cdot (1-L))} \\
  &\leq \norm{\-{proj}_{K_1} (P_v h \cdot L)} + \norm{P_v h \cdot (1-L)},
\end{align*}
where in the last line, we use the fact that $\norm{\-{proj}_{K_1}} \leq 1$.

\begin{proposition} \label{prop:master-form}
For every $i, \!c, \!d$, we define
\begin{align*}
  R_i(\!c,\!d) := \prod_{j\neq i} \tp{1 + \Psi_j(\!c,\!d)}.
\end{align*}
For every $i,\!c$, define the vector $u_{i,\!c}\in\mathbb R^q$ by setting
$u_{i,\!c}(\!d)=0$ when $\nu(v\!d)=0$; otherwise, set
\begin{align*}
  u_{i,\!c}(\!d) := \frac{h(\!d) \nu(v\!d) R_i(\!c,\!d)}{A_{\!d}p_i},
\end{align*}
where we use $A_{\!d}$ to denote the $\!d$-th row of $A$.
Then
  %Suppose $A = \-{diag}(\*a) + \*1\*1^\intercal$ for some vector $\*a\in \mathbb{R}^q$. It holds that
  \begin{align*}
    \norm{\operatorname{proj}_{K_1}(P_v h \cdot L)}^2
    &= \sum_{i, \!c} \nu(v\!c) u_{i,\!c}^\intercal A \tp{ \-{diag}(\nu^{v\!c}_i)  - \nu^{v\!c}_i(\nu^{v\!c}_i)^{\intercal}} A u_{i,\!c}.
  \end{align*}
\end{proposition}
The proof of \Cref{prop:master-form} is a tedious but straightforward calculation from the definitions and is deferred to \Cref{sec:master-form}.

\begin{proposition}\label{cor:related-to-contraction}
Suppose \Cref{cond:alpha-eps-Delta-contraction} holds.
Then, for $\Delta\geq2\alpha$,
%uniformly over $i\in[d]$ and
%$\!d\in[q]$ with $\nu(v\!d)>0$,
%\begin{align}
%  \frac{1}{(A_{\!d}p_i)^2}
%  \sum_{\!c}
%  \nu(v\!c)\nu^{v\!c}(i\!d)
%  R_i(\!c,\!d)^2
%  \leq
%  \tp{1+O_{\alpha,\kappa}\tp{\frac1d}}p_i(\!d).
%  \label{eq:row-bounded-replica-condition}
%\end{align}
every $h\in K_0$ satisfies
\begin{align}
  \norm{\operatorname{proj}_{K_1}(P_vhL)}^2
  &\leq
  \tp{1+O_{\alpha}\tp{\frac1\Delta}}
  \cdot \frac{1-\epsilon}{\Delta}\cdot 
  \norm{h}^2.
  \label{eq:row-bounded-B-bound}
  %\\
  %&\leq
  %\tp{1+O_{\alpha,\kappa}\tp{\frac1d}}
  %\kappa
  %\max_{\substack{\!d\in[q]\\\nu(v\!d)>0}}
  %\left\{
  %  \nu(v\!d)
  %  \sum_{i=1}^d(Bp_i)(\!d)
  %\right\}
  %\norm{h}^2.
  %\label{eq:row-bounded-B-scalar-bound}
\end{align}
%In the diagonal case $B=\-{diag}(\*a)$, the squared operator norm in
%\eqref{eq:row-bounded-B-bound} is exactly
%\begin{align*}
%  \max_{\substack{\!d\in[q]\\\nu(v\!d)>0}}
%  \left\{
%    \nu(v\!d)\*a_{\!d}^2
%    \sum_{i=1}^d p_i(\!d)
%  \right\}.
%\end{align*}
%In particular, if $\alpha,\kappa=O_\delta(1)$ and $d\to\infty$, then
%\eqref{eq:averaged-replica-condition} holds.
\end{proposition}

\begin{proposition}\label{prop:row-bounded-normalizer}
  Suppose $1\leq d\leq\Delta$,
  \Cref{cond:alpha-eps-Delta-contraction} holds, and
  $\Delta\geq2\alpha$.  Then every $h\in K_0$ satisfies
\begin{align}
  \norm{P_vh(1-L)}^2
  \leq
  O_{\alpha}\tp{
    \frac{\log(2d)}{\Delta^2}
  }
  \norm{h}^2.
  \label{eq:row-bounded-normalizer-bound}
\end{align}
\end{proposition}
The proofs of \Cref{cor:related-to-contraction} and \Cref{prop:row-bounded-normalizer} are deferred to
\Cref{sec:related-to-contraction} and \Cref{sec:row-bounded-normalizer}, respectively.
%
Now, we are ready to prove \eqref{eq:general-P10}.

\begin{proof}[Proof of \eqref{eq:general-P10}]
  Combine \Cref{cor:related-to-contraction} and \Cref{prop:row-bounded-normalizer}.
\end{proof}


\subsection{Proof of \Cref{cor:related-to-contraction}}
\label{sec:related-to-contraction}

\begin{proof}
  Fix $h\in K_0$. Since $A=\*1\*1^\intercal-B$ and
  \[\tp{\mathrm{diag}(\nu^{v\!c}_i) - \nu^{v\!c}_i(\nu^{v\!c}_i)^\intercal}\*1=\*0,\]
 \Cref{prop:master-form} gives
\begin{align*}
  \norm{\operatorname{proj}_{K_1}(P_vhL)}^2
  &=
  \sum_{i,\!c}\nu(v\!c)
  u_{i,\!c}^\intercal B \tp{\mathrm{diag}(\nu^{v\!c}_i) - \nu^{v\!c}_i(\nu^{v\!c}_i)^\intercal} B u_{i,\!c} \\
  &\leq 
  \sum_{i,\!c}\nu(v\!c)
  \abs{u_{i,\!c}^\intercal} B\, \mathrm{diag}(\nu^{v\!c}_i) B \abs{u_{i,\!c}} \\
  &\leq 
  \tp{1 + O_\alpha(\Delta^{-1})} \sum_{i,\!c}\nu(v\!c)
  \abs{u_{i,\!c}^\intercal} B\, \mathrm{diag}(p_i) B \abs{u_{i,\!c}},
\end{align*}
where in both inequalities, we use the fact that $B$ is nonnegative and
\begin{align*}
  \nu^{v\!c}_i(\!a) &= \frac{A_{\!c,\!a}p_i(\!a)}{(Ap_i)(\!c)} 
  = \frac{A_{\!c,\!a} p_i(\!a)}{1-(Bp_i)(\!c)}
  \leq \frac{p_i(\!a)}{1-\alpha/\Delta}.
\end{align*}
The last inequality follows from $A \leq 1$ and $(Bp_i)(\!c) \leq \alpha/\Delta$.
Plugging in the definition of $u_{i,\!c}$ from \Cref{prop:master-form}, we have
\begin{align*}
  &\norm{\operatorname{proj}_{K_1}(P_vhL)}^2 \\
  &\leq \tp{1 + O_\alpha(\Delta^{-1})} \sum_{i,\!c} \nu(v\!c) \sum_{\!d,\!a} \abs{u_{i,\!c}(\!d)} (B\mathrm{diag}(p_i)B)_{\!d,\!a} \abs{u_{i,\!c}(\!a)} \\
  &= \tp{1 + O_\alpha(\Delta^{-1})} \sum_{i,\!c} \nu(v\!c) \sum_{\!d,\!a} \frac{\abs{h(\!d)\nu(v\!d)}\abs{h(\!a)\nu(v\!a)}R_i(\!c,\!d)R_i(\!c,\!a)}{(Ap_i)(\!d)(Ap_i)(\!a)} (B\mathrm{diag}(p_i)B)_{\!d,\!a} \\
  &= \tp{1 + O_\alpha(\Delta^{-1})} \sum_{i} \sum_{\!d,\!a} \frac{\abs{h(\!d)\nu(v\!d)}\abs{h(\!a)\nu(v\!a)}(B\mathrm{diag}(p_i)B)_{\!d,\!a}}{(Ap_i)(\!d)(Ap_i)(\!a)} \sum_{\!c}\nu(v\!c) R_i(\!c,\!d)R_i(\!c,\!a).
  %&\leq \tp{1 + O_\alpha(\Delta^{-1})} \sum_{i} \sum_{\!d,\!a} \frac{\abs{h(\!d)\nu(v\!d)}\abs{h(\!a)\nu(v\!a)}(B\mathrm{diag}(p_i)B)_{\!d,\!a}}{(Ap_i)(\!d)(Ap_i)(\!a)} \sum_{\!c}\nu(v\!c) R_i(\!c,\!d)R_i(\!c,\!a) \\
\end{align*}
We claim that for every $\!d$ and $i\in [d]$,
\begin{align} \label{eq:row-bounded-averaged-R-target}
  \sum_{\!c}\nu(v\!c)R_i(\!c,\!d)^2 \leq 1+O_\alpha(\Delta^{-1}).
\end{align}
By the Cauchy--Schwarz inequality, this implies that
\begin{align*}
  \sum_{\!c}\nu(v\!c) R_i(\!c,\!d)R_i(\!c,\!a)
  &\leq  \tp{
    \sum_{\!c}\nu(v\!c)R_i(\!c,\!d)^2
  }^{1/2}
  \tp{
    \sum_{\!c}\nu(v\!c)R_i(\!c,\!a)^2
  }^{1/2}
  \leq 1+O_\alpha(\Delta^{-1}).
\end{align*}
Together, we have
\begin{align*}
  \norm{\operatorname{proj}_{K_1}(P_vhL)}^2
  &\leq \tp{1 + O_\alpha(\Delta^{-1})} \sum_{i} \sum_{\!d,\!a} \frac{\abs{h(\!d)\nu(v\!d)}\abs{h(\!a)\nu(v\!a)}(B\mathrm{diag}(p_i)B)_{\!d,\!a}}{(Ap_i)(\!d)(Ap_i)(\!a)} \\
  &= \tp{1 + O_\alpha(\Delta^{-1})} \sum_{i} \sum_{\!d,\!a} \frac{\abs{h(\!d)\nu(v\!d)}\abs{h(\!a)\nu(v\!a)}(B\mathrm{diag}(p_i)B)_{\!d,\!a}}{(1-(Bp_i)(\!d))(1-(Bp_i)(\!a))} \\
  &\leq \tp{1 + O_\alpha(\Delta^{-1})} \sum_{i} \sum_{\!d,\!a} \abs{h(\!d)\nu(v\!d)}\abs{h(\!a)\nu(v\!a)}(B\mathrm{diag}(p_i)B)_{\!d,\!a} \\
  &= \tp{1 + O_\alpha(\Delta^{-1})} \sum_{i} \abs{h}^\intercal \mathrm{diag}(\nu_v) B\mathrm{diag}(p_i)B \mathrm{diag}(\nu_v) \abs{h} \\
  &= \tp{1 + O_\alpha(\Delta^{-1})} \abs{h}^\intercal \mathrm{diag}(\nu_v) B\mathrm{diag}\tp{\textstyle \sum_i p_i}B \mathrm{diag}(\nu_v) \abs{h} \\
  &\leq \tp{1 + O_\alpha(\Delta^{-1})} \norm{ \-{diag}\tp{\sum_{i=1}^d p_i}^{1/2} B \-{diag}\tp{\nu_v}^{1/2}}_2^2 \norm{h}^2 \\
  \tag{$\epsilon$-contraction}
  &\leq \tp{1 + O_\alpha(\Delta^{-1})} \frac{1-\epsilon}{\Delta}\cdot \norm{h}^2.
\end{align*}
This finishes the proof of \Cref{cor:related-to-contraction}.
%
It only remains to prove \eqref{eq:row-bounded-averaged-R-target}.
By definition, %we have
\begin{align*}
  R_i(\!c,\!d)^2
  &\leq \prod_j\tp{1 + \abs{\Psi_j}(\!c,\!d)}^2 
  \leq \exp\tp{\sum_j 2\abs{\Psi_j}(\!c,\!d)} \\
  \tag{$\e^x \leq 1 + x\e^x$}
  &\leq 1 + \exp\tp{\sum_j 2\abs{\Psi_j}(\!c,\!d)} \cdot \tp{\sum_j 2\abs{\Psi_j}(\!c,\!d)}.
\end{align*}
Hence
\begin{align*}
  \sum_{\!c}\nu(v\!c)R_i(\!c,\!d)^2
  &\leq 1 + O_\alpha(1)\sum_j\sum_{\!c}\nu(v\!c)\abs{\Psi_j}(\!c,\!d) \\
  &=1+O_\alpha(1)\sum_j(\abs{\Psi_j}\nu_v)(\!d) \\
  \tag{\Cref{cor:Psi-related-bounds}}
  &= 1 + O_\alpha(d/\Delta^2)
  = 1 + O_\alpha(\Delta^{-1}).
\end{align*}
This finishes the proof.
%where the last equation follows from \Cref{cor:Psi-related-bounds}.
\end{proof}  


\subsection{Proof of \Cref{prop:row-bounded-normalizer}}
\label{sec:row-bounded-normalizer}

\begin{proof}
Fix $h\in K_0$.  Since $h$ depends only on the center, we write its
value at color $\!c$ as $h(\!c)$.  Center colors of zero $\nu$-mass are
omitted below.  Set
\begin{align*}
  \rho:=\bigotimes_{i=1}^d p_i.
\end{align*}
By \eqref{eq:def-local-mu}, \eqref{def-Pi}, and \Cref{def:L}, the
conditional leaf law $\nu^{v\!c}$ is absolutely continuous with
respect to $\rho$.  Hence, for every $x\in\-{supp}(\rho)$,
\begin{align}
  L(x)
  &=
  \sum_{\!c}\nu(v\!c)\frac{\nu^{v\!c}(x)}{\rho(x)},
  &
  P_vh(x)L(x)
  &=
  \sum_{\!c}\nu(v\!c)h(\!c)
  \frac{\nu^{v\!c}(x)}{\rho(x)}.
  \label{eq:normalizer-likelihood-identities}
\end{align}
Changing measure from the leaf marginal of $\nu$ to $\rho$ gives the
exact identity
\begin{align}
  \norm{P_vh(1-L)}^2
  =
  \E[\rho]{
    (P_vhL)^2\frac{(1-L)^2}{L}
  }.
  \label{eq:normalizer-exact-error}
\end{align}
The integrand on the right is defined to be zero when $L=0$; this is
consistent because \eqref{eq:normalizer-likelihood-identities} then
gives $P_vhL=0$.

It remains to estimate the right-hand side of
\eqref{eq:normalizer-exact-error}.  Values of $\Delta$ bounded in terms
of $\alpha$ can be absorbed into the implicit constant in
\eqref{eq:row-bounded-normalizer-bound}: indeed,
\eqref{eq:normalizer-uniform-bound} and the contraction of $P_v$ give
\begin{align*}
  \norm{P_vh(1-L)}
  \leq
  \tp{1+\exp(2\alpha)}\norm{h}.
\end{align*}
We may therefore assume that $\Delta$ is sufficiently large in terms
of $\alpha$.  Choose
\begin{align*}
  t
  :=
  C_\alpha
  \sqrt{
    \frac1\Delta
    +\frac{d\log(2d)}{\Delta^2}
  },
\end{align*}
where the constant is sufficiently large.  Uniformly over
$1\leq d\leq\Delta$, we then have $t\leq1/2$, and $t$ satisfies the
hypothesis of \Cref{prop:normalizer-conditional-tail}.  Moreover,
\begin{align*}
  \frac{\Delta^2}{d}t^2
  =
  C_\alpha^2
  \tp{
    \frac{\Delta}{d}+\log(2d)
  }
  \geq
  C_\alpha^2\log(2\Delta),
\end{align*}
where the last inequality uses
$\Delta/d\geq\log(\Delta/d)$.  After increasing $C_\alpha$ if
necessary, \Cref{prop:normalizer-conditional-tail} gives
\begin{align}
  \sup_{\!c:\,\nu(v\!c)>0}
  \Pr[\nu^{v\!c}]{\abs{L-1}>t}
  \leq
  \frac{2}{\Delta^3}.
  \label{eq:normalizer-tail-at-t}
\end{align}

On the event $\abs{L-1}\leq t$, we have $L\geq1/2$.  Hence
\eqref{eq:general-P00-second-moment-target} and
\eqref{eq:normalizer-exact-error} give
\begin{align}
  \E[\rho]{
    \*1_{\set{\abs{L-1}\leq t}}
    (P_vhL)^2\frac{(1-L)^2}{L}
  }
  &\leq
  2t^2\E[\rho]{(P_vhL)^2}
  \notag\\
  &\leq
  C_\alpha
  \frac{\log(2d)}{\Delta^2}\norm{h}^2.
  \label{eq:normalizer-good-event}
\end{align}
For the complementary event, weighted Cauchy--Schwarz and
\eqref{eq:normalizer-likelihood-identities} give, for every
$x\in\-{supp}(\rho)$,
\begin{align*}
  (P_vh(x)L(x))^2
  &=
  \tp{
    \sum_{\!c}\nu(v\!c)h(\!c)
    \frac{\nu^{v\!c}(x)}{\rho(x)}
  }^2
  \\
  &\leq
  \tp{
    \sum_{\!c}\nu(v\!c)
    \frac{\nu^{v\!c}(x)}{\rho(x)}
  }
  \tp{
    \sum_{\!c}\nu(v\!c)h(\!c)^2
    \frac{\nu^{v\!c}(x)}{\rho(x)}
  }
  \\
  &=
  L(x)
  \sum_{\!c}\nu(v\!c)h(\!c)^2
  \frac{\nu^{v\!c}(x)}{\rho(x)}.
\end{align*}
Together with \eqref{eq:normalizer-uniform-bound} and
\eqref{eq:normalizer-tail-at-t}, this yields
\begin{align}
  &\E[\rho]{
    \*1_{\set{\abs{L-1}>t}}
    (P_vhL)^2\frac{(1-L)^2}{L}
  }
  \notag\\
  &\qquad\leq
  C_\alpha
  \sum_{x\in\-{supp}(\rho)}
  \rho(x)\*1_{\set{\abs{L(x)-1}>t}}
  \sum_{\!c}\nu(v\!c)h(\!c)^2
  \frac{\nu^{v\!c}(x)}{\rho(x)}
  \notag\\
  &\qquad=
  C_\alpha
  \sum_{\!c}\nu(v\!c)h(\!c)^2
  \sum_{x\in\-{supp}(\rho)}
  \nu^{v\!c}(x)\*1_{\set{\abs{L(x)-1}>t}}
  \notag\\
  &\qquad=
  C_\alpha
  \sum_{\!c}\nu(v\!c)h(\!c)^2
  \Pr[\nu^{v\!c}]{\abs{L-1}>t}
  \notag\\
  &\qquad\leq
  \frac{C_\alpha}{\Delta^3}\norm{h}^2
  \leq
  C_\alpha\frac{\log(2d)}{\Delta^2}\norm{h}^2.
  \label{eq:normalizer-bad-event}
\end{align}
Combining \eqref{eq:normalizer-good-event} and
\eqref{eq:normalizer-bad-event} in
\eqref{eq:normalizer-exact-error} proves
\eqref{eq:row-bounded-normalizer-bound}.
\end{proof}



\subsection{Proof of \Cref{prop:master-form}}
\label{sec:master-form}

\begin{proof}
Fix $h\in K_0$.  In this proposition, $L$ denotes $L_\rho$ for the
reference product distribution
\begin{align*}
  \rho:=\bigotimes_{j=1}^d p_j.
\end{align*}
By \eqref{eq:def-local-mu}, the leaf marginal of $\nu$ is absolutely
continuous with respect to $\rho$.  Thus every sum containing a quotient
by $p_j(\!e)$ may be restricted to $\!e\in\-{supp}(p_j)$; we use this
convention throughout the proof.  Center colors of zero $\nu$-mass do
not contribute to any of the sums below and are omitted.  The
conditional leaf laws on such center fibers may be chosen arbitrarily,
and for every $\!d$ with $\nu(v\!d)=0$, we take
$u_{i,\!c}(\!d)=0$.  For $\nu(v\!d)>0$, \eqref{eq:def-local-mu}
implies that $A_{\!d}p_j>0$ for every $j$, so all the ratios used below
are well-defined.

We first verify that $h$ depends only on the center.  Since
$h\in\Ker(\+L_{N(v)})$, \eqref{def-LB} and the fact that every
$\+L_i$ is an orthogonal projection give
\begin{align*}
  0
  =\inner{h}{\+L_{N(v)}h}
  =\sum_{i=1}^d\inner{h}{\+L_i h}
  =\sum_{i=1}^d\norm{\+L_i h}^2.
\end{align*}
Hence $\+L_i h=0$, or equivalently $P_i h=h$, for every $i$.
Conditional independence in \eqref{eq:def-local-mu} then yields
\begin{align*}
  h=\prod_{i=1}^dP_i h=\E[\nu]{h\mid X_v}.
\end{align*}
We may therefore write $h(\!d)$ for the value of $h$ when
$X_v=\!d$.

Set $F:=P_vhL$.  By \eqref{def-Pi}, both $P_vh$ and $L$ depend only on
the leaf configuration, so the same is true of $F$.  The function $F$
need not be centered, so we apply
\Cref{def:Hoeffding-expand,prop:Hoeffding-expand} to
$F-\E[\nu]{F}$.  Since $K_1\subseteq\*1^\perp$, subtracting this
constant does not change the projection onto $K_1$.  For each
$i\in[d]$, its degree-one component is
\begin{align*}
  \+L_i\prod_{j\ne i}P_jF
  &=
  \E[\nu]{F\mid X_v,X_i}-\E[\nu]{F\mid X_v}.
\end{align*}
Here the constant subtracted from $F$ disappears because
$P_j\*1=\*1$ and $\+L_i\*1=\*0$.
Indeed, conditional independence of the leaves shows that the product
of the projections with $j\ne i$ averages out precisely those leaves,
and applying $P_i$ then averages out the last leaf.  Orthogonality of
the degree-one components now gives
\begin{align}
  \norm{\operatorname{proj}_{K_1}F}^2
  &=
  \sum_{i=1}^d
  \E[\nu]{
    \tp{
      \E[\nu]{F\mid X_v,X_i}
      -\E[\nu]{F\mid X_v}
    }^2
  }
  \notag\\
  &=
  \sum_{i=1}^d\sum_{\!c}
  \nu(v\!c)
  \Var[\nu^{v\!c}_i]{
    \E[\nu]{F(X)\mid X_v=\!c,X_i}
  }.
  \label{eq:master-fiber-variance}
\end{align}

By Bayes' rule, we have
\begin{align*}
  %F(x_{\sim i}, x_i) =(Aw_{i,x_{\sim i}})(x_i), \quad \text{where} \quad
  %w_{i,x_{\sim i}}(\!c)
  %:=
  %\frac{\nu_v(\!c)h(\!c)}{(Ap_i)(\!c)}
  %\prod_{j\ne i}L^{\!c}_j,\quad \forall \!c.
  F(x) = \sum_{\!d}\nu_v(\!d)h(\!d)L^{\!d}(x),
\end{align*}
which implies that
\begin{align*}
   \E[\nu]{F(X)\mid X_v=\!c,X_i=\!a}
   &= \sum_{\!d}\nu(v\!d)h(\!d) \E[x\sim \nu^{v\!c}]{L^{\!d}_i(\!a) \prod_{j\neq i} L^{\!d}_j(x_j)} \\
  \tag{independence}
   &= \sum_{\!d}\nu(v\!d)h(\!d) L^{\!d}_i(\!a) \prod_{j\neq i} \E[\nu^{v\!c}_j]{L^{\!d}_j} \\
   &= \sum_{\!d}\nu(v\!d)h(\!d) L^{\!d}_i(\!a) \prod_{j\neq i} \E[p_j]{L^{\!c}_j L^{\!d}_j} \\
   &= \sum_{\!d}\nu(v\!d)h(\!d) L^{\!d}_i(\!a) \prod_{j\neq i} \tp{1 + \Psi_j(\!c,\!d)} \\
   &= \sum_{\!d}\nu(v\!d)h(\!d) \frac{A_{\!a,\!d}}{(Ap_i)(\!d)} R_i(\!c,\!d) = (Au_{i,\!c})(\!a).
\end{align*}
We finish the proof by applying \Cref{fact:var-matrix-form} to $\Var[\nu^{v\!c}_i]{g}$, where
\begin{align*}
  g(\!a) &:= \E[\nu]{F(X)\mid X_v=\!c,X_i=\!a}. \qedhere
\end{align*}
%
%Hence, we have
%\begin{align*}
%  \E[\nu]{F(X) \mid X_v=\!c,X_i = \!a}
%  &= \E[x \sim \mu^{v\!c}_{N(v)}]{(A w_{i,x_{\sim i}})(\!a)} \\
%  &= A_{\!a} \E[x \sim \mu^{v\!c}_{N(v)}]{w_{i,x_{\sim i}}}
%\end{align*}
%
%
%It remains to compute the conditional expectation in
%\eqref{eq:master-fiber-variance}.  Let
%$x=(x_1,\dots,x_d)$ be in the support of the leaf marginal of $\nu$,
%and write $\nu(x)=\sum_{\!d}\nu(v\!d,x)$ for this marginal.
%By \eqref{def-Pi}, the definition of $L_\rho$, and Bayes' rule,
%\begin{align}
%  F(x)
%  &=
%  L(x)\sum_{\!d}h(\!d)
%  \Pr[\nu]{X_v=\!d\mid X_{[d]}=x}
%  \notag\\
%  &=
%  \frac{\nu(x)}{\rho(x)}
%  \sum_{\!d}h(\!d)\frac{\nu(v\!d,x)}{\nu(x)}
%  \notag\\
%  &=
%  \sum_{\!d}h(\!d)\frac{\nu(v\!d,x)}{\rho(x)}
%  \notag\\
%  &=
%  \sum_{\!d}h(\!d)\nu(v\!d)
%  \prod_{j=1}^d
%  \frac{\nu^{v\!d}(j x_j)}{p_j(x_j)}.
%  \label{eq:master-density-expansion}
%\end{align}
%The final equality follows from the conditional product structure in
%\eqref{eq:def-local-mu} and from
%$\rho(x)=\prod_{j=1}^dp_j(x_j)$.
%
%Fix $i\in[d]$, a center color $\!c$ with $\nu(v\!c)>0$, and a leaf
%color $\!a\in\-{supp}(\nu^{v\!c}_i)$.  Let $\sigma$ range over the
%configurations of the leaves other than $i$, and let
%$(\sigma;i\!a)$ denote the full leaf configuration obtained by placing
%$\!a$ at coordinate $i$.  Conditional independence and
%\eqref{eq:master-density-expansion} give
%\begin{align*}
%  &\E[\nu]{F(X)\mid X_v=\!c,X_i=\!a}
%  \\
%  &=
%  \sum_{\sigma}
%  \prod_{j\ne i}\nu^{v\!c}(j\sigma_j)
%  F(\sigma;i\!a)
%  \\
%  &=
%  \sum_{\sigma}
%  \prod_{j\ne i}\nu^{v\!c}(j\sigma_j)
%  \sum_{\!d}h(\!d)\nu(v\!d)
%  \frac{\nu^{v\!d}(i\!a)}{p_i(\!a)}
%  \prod_{j\ne i}
%  \frac{\nu^{v\!d}(j\sigma_j)}{p_j(\sigma_j)}
%  \\
%  &=
%  \sum_{\!d}h(\!d)\nu(v\!d)
%  \frac{\nu^{v\!d}(i\!a)}{p_i(\!a)}
%  \sum_{\sigma}
%  \prod_{j\ne i}
%  \frac{
%    \nu^{v\!c}(j\sigma_j)\nu^{v\!d}(j\sigma_j)
%  }{p_j(\sigma_j)}
%  \\
%  &=
%  \sum_{\!d}h(\!d)\nu(v\!d)
%  \frac{\nu^{v\!d}(i\!a)}{p_i(\!a)}
%  \prod_{j\ne i}
%  \tp{
%    \sum_{\!e\in\-{supp}(p_j)}
%    \frac{\nu^{v\!c}(j\!e)\nu^{v\!d}(j\!e)}{p_j(\!e)}
%  }
%  \\
%  &=
%  \sum_{\!d}h(\!d)\nu(v\!d)R_i(\!c,\!d)
%  \frac{\nu^{v\!d}(i\!a)}{p_i(\!a)}.
%\end{align*}
%The finite sum over $\sigma$ factors because its summand is a product
%over $j\ne i$.
%
%The one-leaf posterior law in \eqref{eq:def-local-mu} gives
%\begin{align*}
%  \frac{\nu^{v\!d}(i\!a)}{p_i(\!a)}
%  =
%  \frac{A_{\!d,\!a}}{A_{\!d}p_i}.
%\end{align*}
%Therefore, by the definition of $u_{i,\!c}$ and the symmetry of $A$,
%\begin{align}
%  \E[\nu]{F(X)\mid X_v=\!c,X_i=\!a}
%  &=
%  \sum_{\!d}u_{i,\!c}(\!d)A_{\!d,\!a}
%  \notag\\
%  &=
%  (A^\intercal u_{i,\!c})(\!a)
%  =
%  (A u_{i,\!c})(\!a).
%  \label{eq:master-one-leaf-expectation}
%\end{align}
%Outside $\-{supp}(\nu^{v\!c}_i)$, we choose
%$Au_{i,\!c}$ as a version of this conditional expectation.  These
%coordinates do not affect the variance, since the corresponding rows
%and columns of the covariance matrix below vanish.
%
%For brevity, set
%\begin{align*}
%  M_{i,\!c}
%  :=
%  \-{diag}(\nu^{v\!c}_i)
%  -\nu^{v\!c}_i(\nu^{v\!c}_i)^\intercal.
%\end{align*}
%Expanding the scalar variance in \eqref{eq:master-fiber-variance} and
%using \eqref{eq:master-one-leaf-expectation}, we obtain
%\begin{align*}
%  &\Var[\nu^{v\!c}_i]{
%    \E[\nu]{F(X)\mid X_v=\!c,X_i}
%  }
%  \\
%  &=
%  \sum_{\!a}\nu^{v\!c}_i(\!a)(Au_{i,\!c})(\!a)^2
%  -
%  \tp{
%    \sum_{\!a}\nu^{v\!c}_i(\!a)(Au_{i,\!c})(\!a)
%  }^2
%  \\
%  &=
%  (Au_{i,\!c})^\intercal M_{i,\!c}(Au_{i,\!c}).
%\end{align*}
%Substitution into \eqref{eq:master-fiber-variance} yields
%\begin{align*}
%  \norm{\operatorname{proj}_{K_1}F}^2
%  &=
%  \sum_{i,\!c}\nu(v\!c)
%  (A u_{i,\!c})^\intercal
%  M_{i,\!c}
%  (A u_{i,\!c}). \qedhere
%  %\label{eq:master-before-diag}
%\end{align*}
%
%Finally, since $\nu^{v\!c}_i$ is a probability vector,
%$M_{i,\!c}\*1=\*0$.  The matrix $M_{i,\!c}$ is symmetric, so also
%$\*1^\intercal M_{i,\!c}=\*0$.  The identity
%$A=\-{diag}(\*a)+\*1\*1^\intercal$ thus gives
%\begin{align*}
%  (Au_{i,\!c})^\intercal M_{i,\!c}(Au_{i,\!c})
%  &=
%  u_{i,\!c}^\intercal
%  (A^\intercal-\*1\*1^\intercal)
%  M_{i,\!c}
%  (A-\*1\*1^\intercal)
%  u_{i,\!c}
%  \\
%  &=
%  u_{i,\!c}^\intercal
%  \-{diag}(\*a)M_{i,\!c}\-{diag}(\*a)
%  u_{i,\!c}.
%\end{align*}
%Substituting this identity into \eqref{eq:master-before-diag} proves the
%proposition.
\end{proof}

%\subsection{An averaged replica condition}
%\label{subsec:averaged-replica-condition}

\section{Application to the anti-ferromagnetic Potts model}
\label{sec:Potts}
In this section, we prove \Cref{thm:Potts}.
Let $G=(V,E)$ be an $n$-vertex simple graph with maximum degree
$\Delta$ and girth at least $5$.  Fix $\delta\in(0,1)$,
$\beta\in[0,1]$, and an integer $q\geq2$ satisfying
\begin{align*}
  q\geq(1+\delta)(1-\beta)\Delta.
\end{align*}
Recall that the anti-ferromagnetic $q$-state Potts distribution
$\mu_G$ on $[q]^V$ is defined by
\begin{align*}
  \mu_G(\sigma)
  \propto
  \beta^{\abs{\set{uv\in E:\sigma_u=\sigma_v}}},
  \qquad \sigma\in[q]^V.
\end{align*}
Thus every edge whose endpoints have the same color contributes a
factor $\beta$, and there is no external field.  All conditioned laws
and marginals below are formed from $\mu_G$.

\begin{proof}[Proof of \Cref{thm:Potts}]
We finish the proof by applying \Cref{thm:general}.
To prove the spectral-gap bound, it suffices to verify
\Cref{cond-local-spectral-contraction}.
Put $\vartheta:=1-\beta$.  The interaction matrix and its deficit matrix
are
\begin{align*}
  A=\*1\*1^\intercal-\vartheta\mathbb I,
  \qquad
  B:=\*1\*1^\intercal-A=\vartheta\mathbb I.
\end{align*}
We verify \Cref{cond-local-spectral-contraction} with
\begin{align*}
  \alpha:=\frac1\delta,
  \qquad
  \epsilon:=\frac{\delta}{1+\delta}.
\end{align*}

Fix $v\in V$ and a feasible pinning $\tau$ outside $N[v]$, and put
$d:=\deg_G(v)$.  For $i\in N(v)$ and $\!a\in[q]$, let
$k_{i,\!a}$ be the number of pinned neighbors of $i$ outside $N[v]$
whose color is $\!a$.  Since $G$ has girth at least $5$, $N[v]$
induces a star.  Thus, writing
\begin{align*}
  p_i:=\mu^\tau_{G-v,i},
\end{align*}
we have
\begin{align*}
  p_i(\!a)
  =
  \frac{\beta^{k_{i,\!a}}}
       {\sum_{\!b=1}^q\beta^{k_{i,\!b}}}.
\end{align*}
When $\beta=0$, we use the convention $0^0=1$; feasibility of $\tau$
ensures that the denominator is positive.  Bernoulli's inequality and
$\sum_{\!b}k_{i,\!b}\leq\Delta-1$ give
\begin{align*}
  \sum_{\!b=1}^q\beta^{k_{i,\!b}}
  &=
  \sum_{\!b=1}^q(1-\vartheta)^{k_{i,\!b}}
  \geq
  q-\vartheta\sum_{\!b=1}^q k_{i,\!b}
  \geq
  q-\vartheta(\Delta-1).
\end{align*}
Consequently, if $\vartheta>0$, then the assumption on $q$ yields
\begin{align}
  \vartheta p_i(\!a)
  \leq
  \frac{\vartheta}{q-\vartheta(\Delta-1)}
  \leq
  \frac1{\delta\Delta+1}
  \leq
  \frac1{\delta\Delta}.
  \label{eq:Potts-leaf-interaction-bound}
\end{align}
The same conclusion is immediate when $\vartheta=0$.

For each $\!a\in[q]$, set
\begin{align*}
  w_{\!a}
  &:=%
  \prod_{i\in N(v)}\tp{1-\vartheta p_i(\!a)},
  &
  Z&:=\sum_{\!b=1}^q w_{\!b}.
\end{align*}
The tree recursion on the conditioned star gives
\begin{align*}
  \nu_v(\!a)
  :=
  \mu^\tau_{G,v}(\!a)
  =
  \frac{w_{\!a}}Z.
\end{align*}
Since $\prod_i(1-u_i)\geq1-\sum_i u_i$ for $u_i\in[0,1]$,
\begin{align*}
  Z
  \geq
  \sum_{\!a=1}^q\tp{1-\vartheta\sum_{i\in N(v)}p_i(\!a)}
  =q-\vartheta d.
\end{align*}
It follows, again trivially when $\vartheta=0$, that
\begin{align}
  \vartheta\nu_v(\!a)
  \leq
  \frac{\vartheta}{q-\vartheta d}
  \leq
  \frac1{\delta\Delta}.
  \label{eq:Potts-center-interaction-bound}
\end{align}
Because $B=\vartheta\mathbb I$, \eqref{eq:Potts-leaf-interaction-bound}
and \eqref{eq:Potts-center-interaction-bound} prove the
$\alpha$-bounded marginal condition.

It remains to verify the contraction condition.  The case
$\vartheta=0$ is immediate, so assume $\vartheta>0$.  For each
$\!a\in[q]$, put
\begin{align*}
  T_{\!a}:=\vartheta\sum_{i\in N(v)}p_i(\!a).
\end{align*}
By \eqref{eq:Potts-leaf-interaction-bound} and
$\log(1-u)\geq-u/(1-u)$ for $u\in[0,1)$,
\begin{align*}
  \log w_{\!a}
  &\geq
  -\tp{1+\frac1{\delta\Delta}}T_{\!a}.
\end{align*}
Since $\sum_{\!a}T_{\!a}=\vartheta d$, the arithmetic--geometric
mean inequality gives
\begin{align}
  Z
  \geq
  q\tp{\prod_{\!a=1}^q w_{\!a}}^{1/q}
  \geq
  q\exp\tp{
    -\tp{1+\frac1{\delta\Delta}}
    \frac{\vartheta d}{q}
  }.
  \label{eq:Potts-normalizer-lower-bound}
\end{align}
On the other hand, $w_{\!a}\leq\exp(-T_{\!a})$.  Therefore
$T_{\!a}\exp(-T_{\!a})\leq \e^{-1}$ and
\eqref{eq:Potts-normalizer-lower-bound} imply
\begin{align*}
  &\Delta\vartheta^2\nu_v(\!a)
    \sum_{i\in N(v)}p_i(\!a)
  \\
  &\qquad=
  \Delta\vartheta T_{\!a}\frac{w_{\!a}}Z
  \\
  &\qquad\leq
  \frac{\vartheta\Delta}{q}
  \exp\tp{
    \tp{1+\frac1{\delta\Delta}}
    \frac{\vartheta d}{q}-1
  }.
\end{align*}
For sufficiently large $\Delta$ in terms of $\delta$, we may assume
$\Delta\geq\delta^{-2}$.  Using $d\leq\Delta$ and
$\vartheta\Delta/q\leq1/(1+\delta)$, we then have
\begin{align*}
  \tp{1+\frac1{\delta\Delta}}
  \frac{\vartheta d}{q}
  \leq1.
\end{align*}
Hence, uniformly in $\!a$,
\begin{align*}
  \vartheta^2\nu_v(\!a)
  \sum_{i\in N(v)}p_i(\!a)
  \leq
  \frac1{(1+\delta)\Delta}
  =
  \frac{1-\epsilon}{\Delta}.
\end{align*}
Since $B=\vartheta\mathbb I$, the matrix in the contraction condition
is diagonal and therefore
\begin{align*}
  &\norm{
    \-{diag}\tp{\sum_{i\in N(v)}p_i}^{1/2}
    B
    \-{diag}(\nu_v)^{1/2}
  }_2^2
  \\
  &\qquad=
  \vartheta^2\max_{\!a\in[q]}\set{
    \nu_v(\!a)\sum_{i\in N(v)}p_i(\!a)
  }
  \leq
  \frac{1-\epsilon}{\Delta}.
\end{align*}
The estimates above are uniform over $v$ and all feasible pinnings
$\tau$.  Thus the Potts measure has local
$(\delta^{-1},\delta/(1+\delta))$-spectral contraction.

The Glauber dynamics is irreducible.  Indeed, when $\beta>0$ the Gibbs
measure has full support.  When $\beta=0$, the assumption on $q$ gives
$q\geq\Delta+2$ for sufficiently large $\Delta$ in terms of $\delta$,
and the proper-coloring Glauber dynamics is irreducible in this range.
After increasing the lower bound on $\Delta$ once more if necessary,
all the hypotheses of \Cref{thm:general} hold.  We conclude that the
discrete-time Glauber dynamics has spectral gap
$\Omega_\delta(1/n)$.

Finally, consider the smallest positive mass of a configuration.  If
$\beta=0$, the measure is uniform on a set of at most $q^n$ proper
colorings, so
\begin{align*}
  \min_{\sigma\in\-{supp}(\mu)}\mu(\sigma)\geq q^{-n}.
\end{align*}
If $0<\beta\leq1$, then every configuration has unnormalized weight at
least $\beta^{\abs{E}}\geq\beta^{n\Delta/2}$, while the partition
function is at most $q^n$.  Hence
\begin{align*}
  \min_{\sigma\in[q]^V}\mu(\sigma)
  \geq
  \beta^{n\Delta/2}q^{-n}.
\end{align*}
Substituting these bounds and the spectral-gap estimate into the
standard worst-start mixing bound
\eqref{eq:spgap-implies-mixing-time} gives, respectively,
\begin{align*}
  t_{\mathrm{mix}}(\varepsilon)
  &=
  O_\delta\tp{
    n^2\log q+n\log\frac1\varepsilon
  },
  &&\beta=0,
  \\
  t_{\mathrm{mix}}(\varepsilon)
  &=
  O_\delta\tp{
    n^2\log q+n^2\Delta\log\frac1\beta
    +n\log\frac1\varepsilon
  },
  &&0<\beta\leq1.
\end{align*}
This proves the theorem.
\end{proof}

\bibliographystyle{alpha}
\begingroup
\footnotesize
\raggedright
\bibliography{refs}
\endgroup

\end{document}
